\chapter{環論}
この章では環論の問題を扱う。環論は代数幾何学, 整数論を展開するために基礎となる分野であり, なかなか手強いものである。参考書としては\cite{atiyah}が有名である。なお, スキーム論の対策はこのpdfの専門のほうに述べる. \par
唐突だが, 次の問題に挑戦してみよう。次の問題は, 設問が5つあるが, だんだんレベルが上がっていくような構成が見事な問題である。
\begin{prac} (神戸大R1-A) \label{prac:kobeR1_A}
$p$を素数とする。
\[R = \set{\dfrac{m}{n}\in \mb{Q}\mid m,n\in \mb{Z}, p\not | n}\]
とおく。次の問いに答えよ。
\ben
\item $R$が$\mb{Q}$の部分環であることを示せ。
\item $p$が生成する$R$のイデアルが極大であることを示せ。
\item (2)の極大イデアル以外に$R$が極大イデアルを持たないことを示せ。
\item $R$係数$1$変数多項式環$R[X]$の$X^2+1$で生成されるイデアルを$I$とおく。$I$が素イデアルであることを示せ。
\item (4)のイデアル$I$を含む$R[X]$の極大イデアルを$J$とおく。$J$が$p$を含むことを示せ。
\een
\end{prac}
(5)は単純に見えるが, 可換環論の知識を用いるとテクニカルに導ける。(4)までなら速く解けて, かつこの(5)ができる人は可換環論に一定の親しみがあると思われる。(4)までできれば基礎力はある。ひとまず(4)までを自力でできるようになるとよい。(5)が分からなかった人は, \cite{atiyah}の5章(整従属)のあたりを復習してみよう。
(5)の答えのみ脚注\footnote{
$R$代数の同型$R[X]/(X^2+1)\cong R[\sqrt{-1}]$がある。包含写像$i:R\to R[i]$は単射整準同型であるから, $i$が誘導する$\spec{i}: \spec{R[i]}\to \spec{R}$は極大イデアルを極大イデアルに移す。なぜなら, $R/i^{-1}(J)\to R[i]/J$は整域の間の単射整準同型で,$R[i]/J$が体であることと$R/i^{-1}(J)$が体であることが同値であるためである。よって$i^{-1}(J) = J\cap R$は$R$の極大イデアルだから$J\cap R = (p)$である。ゆえに$p\in J$を得る。\qed

}に与えておく。\par





重要な環と簡単な性質について復習する。\tf{以降ことわりの無い限り環は単位的可換環を指す。}




\section{さまざまな環}
ここでは, 次の種類の環の性質をまとめる.
\ben
\item 整域, UFD, PID. これらは代数入門くらいのものである.
\item ($\spadesuit$) 整閉整域, Noether環,  Artin環. こちらは代数学I程度のものである.
\item ($\spadesuit\spadesuit$) Dedekind環, 離散付値環, 正則局所環, Cohen-Macaulay環, カテナリー環. これらは専門レベルである.
\een

\subsection{整域}

\tbox{
\begin{prop}
$A$を環とする。$A$が整域であるとは, いかなる$x,y\in A\setminus{\{ 0 \}}$に対しても$xy \neq 0$であるという性質を満たす環のことをいう。整域$A$について, 次が成り立つ.
\ben
\item 整域であることと$(0)$を素イデアルに持つことは同値.
\item 整域の局所化は整域. \tf{整域の部分環も整域.}
\item 整域の素イデアル$(0)$による局所化は体である. これを\tf{商体}という.
\een
\end{prop}
}{title=整域のまとめ}

\begin{eg}
$\mb{Z}[X]$は整域である. $\mb{Z}[X]$の商体は$\mb{Q}(X)$.  $\mb{Z}[X]$の極大イデアル$(p,X)$を考える. このとき$\mb{Z}[X]_{p}\cap \mb{Z}[X]_{X} = \mb{Z}[X]$である.
\end{eg}


\begin{egprob}
$A$を$\mb{Z}$の素イデアル$(2)$による局所化とする.すなわち
\[A = \mb{Z}_{(2)} = \set{\dfrac{a}{b} \mid b\notin 2\mb{Z}, a\in \mb{Z}} \subset \mb{Q}\]
とおく. $A$が整域であることを示せ. また, $A$の素イデアル$\maf{p}$をすべて求めよ. また, $A/\maf{p}$は何か?
\end{egprob}
\begin{egans}
$A$が整域なのは部分環だから. 単純に調べることも可能. \\
素イデアル$\maf{p}$は$2\mb{Z}_{(2)}$のただひとつ. $A/\maf{p} \cong \mb{F}_{2}$. これの証明は, $A\to \mb{F}_{2}$を分子の$\bmod{2}$に送る写像とすれば核が$\maf{p}$になるから.
\end{egans}



\begin{egprob}
$\phi: \mb{Z}_{(2)} \to \mb{F}_{2}\times \mb{Q}$を
\[\phi(\dfrac{a}{b}) = (a\bmod{2}, \dfrac{a}{b})\]
により定める.
\ben
\item $\phi$が環準同型であることを示せ.
\item $\mb{F}_{2}\times \mb{Q}$の素イデアルをすべて求めよ. また, それぞれの素イデアル$P\subset \mb{F}_{2}\times \mb{Q}$に対して, $\phi^{-1}(P)$を求め, 素イデアルを引き戻す写像$P\mapsto \phi^{-1}(P)$が$\set{\mb{F}_{2}\times \mb{Q}の素イデアル}$と$\set{\mb{Z}_{(2)}の素イデアル}$の全単射を与えることを示せ.
\item ($\spadesuit\spadesuit$) $\phi^{\ast}: \spec{\mb{F}_{2}\times \mb{Q}}\to \spec{\mb{Z}_{(2)}}$は\tf{同相写像ではない}ことを示せ.
\een
\end{egprob}

\begin{rem}
上の問題は, $\phi^{\ast}$が全単射だが同相でない例を定めている.
\end{rem}



%-------------------------UFD-----------

\subsection{UFD}

整域$A$の元$p$が生成するイデアル$pA$が素イデアルであるとき, $p$を\textbf{素元}という.

\begin{egprob}
素元は既約元であることを証明せよ.
\end{egprob}
\begin{egans}
素元$p$を取る. $p=ab$と分解したとき, $a,b$のどちらかが単元であることを示せばよい. $ab$は$p$で割れるから, $a,b$のどちらかは$p$で割り切れる. たとえば$a=pa'$とする. このとき, 整域なので$p\neq 0$で割ると$1=a'b$となる. よって$b$は単元であり, 示せた. \qed
\end{egans}


\tbox{
\begin{prop}
整域$A$であって, 次の条件$(\ast)$を満たすとき$A$は\tf{一意分解整域(UFD)}という。
\ben
\item[($\ast$)] $A$の$0$でも単元でもない任意の元$a$に対して, ある素元$p_1,\cdots, p_t$が存在して$a=p_1p_2\cdots p_t$が成り立つ。
\een
一意分解整域$A$に対して次が成り立つ。
\ben
\item 素元分解は本質的に一意である。
\item 多項式環$A[x]$もUFDである。
\item \tf{UFDにおいて既約元は素元である. }
\een
\end{prop}
}{title=UFDまとめ}


\begin{eg}
UFDの剰余環は必ずしもUFDではない. たとえば,
\ben
\item $k[X,Y]/(X^2-Y^3)$ (整閉ですらない)
\item $\mb{F}_{3}[X,Y]/(X^2+Y^2-1)$.
\een
\end{eg}

\begin{egprob}
$\mb{F}_{3}[X,Y]/(X^2+Y^2-1)$は整域であることを示せ.
\end{egprob}
\begin{egans}
$\mb{F}_{3}[X,Y]$で$p=Y-1$は素元. $X^2+Y^2-1$はアイゼンシュタイン多項式なので既約. よって整域.
\end{egans}




\begin{eg} (\tf{円の座標環})
標数2でない代数閉体$k$に対して$k[x,y]/(x^2+y^2-1)$はPIDであることを示す. まず代数閉体なので$i^2=-1$なる$i\in k$が存在する. そこで, $x^2+y^2=(x+iy)(x-iy)$と分解すると, $f(x,y) \mapsto f(\frac{u+v}{2}, \frac{u-v}{2i})$ は次の同型を誘導する.
\bealn{
k[x,y]/(x^2+y^2 - 1) \to k[u,v]/(uv-1)
}
右の環は$k[u,1/u]$という環に同型であり, これは1変数多項式環$k[u]$の$\{ 1,u,\cdots\}$による局所化に同型である. 一般にPIDの局所化はPIDであるため\footnote{Hint: 局所化のイデアルは必ず拡大イデアルである. つまり$S^{-1}I$の形で書ける. ここで$S$は積閉集合, $I$は環$A$のイデアル.}, これもPIDであることが従う. なお, ここで用いた$k$の性質は$i^2=-1$となる元の存在正と標数が$2$でないことであるため, $k=\mb{F}_{5}$でも正しい. しかしこのような$i$がないならPIDになるとは限らない.
\end{eg}
\begin{prac}
$k=\mb{F}_{3}$のときにこの環がUFDですらないことを, 次の等式から示せ.
\[(x-y)(x+y) = 2(x^2+1)\]
\end{prac}








\subsection{PID}


\tbox{
\begin{prop}
整域であって, すべてのイデアルが一つの元で生成されるもののことを\tf{単項イデアル整域(PID)}という。PIDに対して次が成り立つ。
\ben
\item PIDはUFDである。
\item PIDの0でない素イデアルは極大イデアルである。とくに, $k[X]$の極大イデアルは$k$上既約な多項式によって生成される。
\item
\een
\end{prop}
}{title=PIDまとめ}
\prf
(2): (\tf{自分でできるように！}) $(x),(y)$が0でない素イデアルで $(x)\subsetneq (y)$なるものとする.  このとき$x=ya$と表せる. もし$y\in (x)$なら$(x)=(y)$となるのでおかしい. ゆえに$y\notin (x)$であり, これと$x=ya$により$a\in (x)$が従う. よって$a=xb$と書くと, $1=ab$となり, $a,b$は単元なので, $x=ya$から$x,y$は同伴. すると$(x) = (y)$なのでおかしい. よって極大性が示せた.

\qed




\begin{eg}
$\mb{C}[x]$はユークリッド整域なのでPIDでありUFDなので整閉整域である.
\end{eg}

\begin{egprob}
単項イデアルの共通部分はまた単項イデアルか?
\end{egprob}
\begin{egans}
ちゃんと調べていないが, この問題は平坦性と関係があるらしい. \par
$A=k[t^2,t^3]$で$(t^2)\cap (t^3) = (t^4,t^5)$になる. これは単項で書けない. $(t^4)$とは等しくないことに注意！
\end{egans}







\begin{eg}
次はPIDである: $\mb{Z}$, $k[X]$ ($k$は体).
\end{eg}

\begin{rem}
任意のイデアルが単項イデアルだが整域でない例がある. たとえば$A=\mb{C}[X]/(X^2)$を考える. これのイデアルは$(0),(X),(1)$に限られる. 実際, この環はいわゆる局所環で, $(X)$に属さない元はすべて単元であるということに注意すれば, イデアル$I$が$(1)$でも$(0)$でもないなら$X$をh含むしかなく, $I=(X)$とならざるを得ない.
\end{rem}

\begin{prac} (代数民なら必ず答えられるようにしたい)\\
PIDでないUFDの例を2個答えよ.
\end{prac}


\subsection{Euclid整域}
Euclid整域も代数入門レベルの対象であるが, 登場頻度がかなり少ないので短めに述べる.
\begin{prop}

\end{prop}


\begin{egprob}
$\mb{Z}[i], \mb{Z}[\omega]$がEuclid整域であることを示せ.
\end{egprob}






%------------------------整閉--------
\subsection{整閉整域}

\tbox{
\begin{prop}
$A$を環とする。$A$の商体\footnote{積閉集合$S=\{ A\text{の零因子でない元} \}$による局所化$S^{-1}A$のことをいう。}を$K$とする。$A$が整域で, $K$の元で$A$上整なものが$A$の元のみであるとき, $A$は\tf{整閉整域}(または正規環)という。整閉整域$A$にたいして次が成り立つ。
\ben
\item $A$の局所化は整閉整域である.
\item $A$が整閉整域ならば多項式環$A[T]$も整閉整域である.
\een
\end{prop}
}{title=整閉整域まとめ}
\begin{prac}

\end{prac}

\begin{rem}
すなわち,$A$が整閉整域とは,  $a\in K$が$A$係数の方程式
\[x^{n} + a_1x^{n-1} + \cdots + a_{n} = 0\]
の根であるならば必ず$a\in A$であるということを言っている。$A=\mb{Z}$ならば$K=\mb{Q}$であるが 「ある有理数$a$が整数係数のモニック多項式の根であるならば$a$は整数である」は高校数学でも聞く話であり, 有理数の分母と分子の素因数分解を考えることで初頭整数論的に証明できる。つまり$\mb{Z}$は整閉整域であるのだが,実はUFDは必ず整閉整域である。その証明は$\mb{Z}$での証明とほとんど平行に進むであろう。
\end{rem}

\begin{recall} $d\neq 0$が平方因子を持たない整数とする.
$\mb{Q}(\sqrt{d})$の整数環は$d\equiv 1\pmod{4}$で$\mb{Z}[(\sqrt{d}+1)/2]$で$d\equiv 2,3\pmod{4}$なら$\mb{Z}[\sqrt{d}]$.
\end{recall}


\begin{prac}
$\mb{Z}$が整閉整域であることを直接示せ. (これは実は高校数学でやっている.)
\end{prac}

\begin{prac} (仮想面接問題)
整閉整域の局所化は整閉整域ですか?
\end{prac}



\begin{egprob} (\cite{aoyukie})
$A=\mb{C}[X,Y,Z]/(Z^2 - XY)$が正規環であることを示せ。
\end{egprob}
\begin{egans}
$X,Y,Z$の$A$での像を$x,y,z$とする. $\mr{Frac}A =:K  = \mb{C}(x,y,z)$である. $z$は$\mb{C}(x,y)$上代数的なので$K=\mb{C}(x,y)[z]$. よって任意の$K$の元は$f(x,y) + zg(x,y)$と表せる. この元が$A$上整であるとせよ. $A$の自己同型$z\mapsto -z$により共役元$f - zg$も$A$上整であり, トレース$2f$もノルム$f^2-z^2g^2 = f^2-xyg^2$も$A$上整である. $A$は$\mb{C}[x,y]$上整だから$f$も$f^2-xyg^2$も$\mb{C}[x,y]$上整であり, $\mb{C}[x,y]$は整閉だから$f\in\mb{C}[x,y]$かつ$xyg^2\in \mb{C}[x,y]$である. $(xyg)^2 \in (x)\cap (y)$で, $x,y$は素元なので$xyg\in (x)\cap (y)=(xy)$となる. よって$g\in \mb{C}[x,y]$だから$f+zg \in \mb{C}[x,y,z]$となり$A$は整閉である. \qed
\end{egans}


\begin{egprob} (H9 数学II-1)\\
$K$は標数が2でないとし, $F=F(X_1,\dots, X_n)\in K[X_1,\dots, X_n]$は$K[X_1,\dots, X_n]$の元の二乗にはなっていないと仮定する. $R=K[X_1,\dots, X_n, Y]/(Y^2 - F(X_1,\dots, X_n))$を考える.
\ben
\item $R$は整域であることを示せ.
\item $K[X_1,\dots, X_n]$が正規環であることを利用して\footnote{UFDであることから直ちに従う.}, $R$が正規環であるための必要十分条件は$F$が二乗因子を持っていないことであることを示せ.
\een
\end{egprob}

\begin{egans}
(1): $R$の$X_i,Y$の像をそれぞれ$x_i,y$で書く. $R$の元は$f(x_1,\dots, x_n) + g(x_1,\dots, x_n)y$一意にで表せる. このような元$f_1 + g_1y, f_2+g_2y$をとり, その積$f_1f_2 + y(f_1g_2 + g_1f_2) + y^2g_1g_2 = (f_1f_2 - Fg_1g_2) + y(f_1g_2 + g_1f_2) $が0であるとしよう. このとき$f_1f_2 - Fg_1g_2 = 0$かつ$f_1g_2 + g_1f_2 = 0$を得る. よって
\[g_1f_1f_2 - Fg_1^2g_2 = -f_1^2g_2 - Fg_1^2g_2 = 0\]
を得る. \par
もし$g_2\neq 0$なら$f_1^2 + Fg_1^2 = 0\in R$. すなわち$f_1^2 + Fg_1^2 \in k[X_1,\dots, X_n]\cap (Y^2-F(X_1,\dots,X_n)) = (0)$だから$f_1^2 + Fg_1^2 = 0 \in K[X_1,\dots, X_n]$である. $g_1\neq 0$なら$F$の素元分解を考えて$F$が平方元となるのでおかしい. よって$g_1=0$とすると$f_1^2=0$で$f_1=0$になるからこの場合はよい. \par
また, $g_2 = 0$で$g_1\neq 0$なら, 先の$g_1,g_2$を入れ替えた議論が可能である. $g_1=g_2=0$の場合, $f_1f_2 = 0 \in K[X_1,\dots, X_n]\cap (Y^2-F(X_1,\dots, X_n)) = (0)$なので$f_1,f_2$のどちらかは0で, この場合も整域の定義を満たす. よって$R$は整域. \qed \\
(2): \bolm{K(x_1,\dots, x_n)(y) = K(x_1,\dots, x_n)[y]}に注意せよ($\because y$が代数的). したがって$\mr{Frac}R$の任意の元は$ f_1(x_1,\dots, x_n) + f_2(x_1,\dots, x_n)y$と表せる($f_1,f_2\in K(x_1\dots, x_n)$).
$F$が二乗因子を持たないとし, $R$が正規であることを証明する. 自然な単射 $\phi:K[X_1,\dots,X_n]\hookrightarrow R$により部分環とみなす. $y$は整等式$y^2 - F = 0$を満たすので, $\phi$は整準同型である.   $\alpha=f_1 + yf_2$が$R$上整であるとする. $K(x_1,\dots, x_n)$上の共役元は$\beta := f_1 - yf_2$で, $\mr{Tr}(\alpha) = \alpha + \beta = 2f_1(x_1,\dots, x_n)$, $\mr{N}(\alpha) = \alpha\beta = f_1^2 - Ff_2^2$も$R$上整である. $\phi$は整準同型なので, $2f_1$と$f_1^2 - Ff_2^2$は$K[X_1,\dots, X_n] \cong K[x_1,\dots, x_n]$ 上整. これは正規環上整であるから, $2f_1, f_1^2 - Ff_2^2 \in K[x_1,\dots, x_n]$が従う. $K$は標数2でないので$f_1\in K[x_1,\dots, x_n]$である. あとは$f_2\in K[x_1,\dots, x_n]$を示せばよいが, それも正しい. 実際, $f_2 = a_2/b_2$で$a_2,b_2\in K[x_1,\dots, x_n]$は互いに素とすると, $Fa_2^2 = gb_2^2 (g\in K[x_1,\dots,x_n])$という式がある. $b_2$がもしある素元$p$で割れるなら, 互いに素の仮定より$a_2^2$は$p$で割れない. $b_2^2$は$p^2$で割れるので, $F$が$p^2$で割れなければならない.  これは平方因子を持たないことに反する. よって$b_2$はいかなる素元でも割れないので単元であり, $b_2\in K\setminus{\set{0}}$が従うので$f_2 = b_2^{-1}a_2 \in K[x_1,\dots, x_n]$となる. よって$\alpha \in K[x_1,\dots, x_n,y]$だから$R$は正規環である. \par
逆に$F(x_1,\dots, x_n)$が平方因子$p^2$を持つとする. $F/p^2 = F_0(x_1,\dots, x_n)$とおく. $f_2 = 1/p$とすると $f_2y \notin K[x_1,\dots, x_n,y]$であるが, これは
\[(f_2y)^2 - F_0 = 0\]
という整等式を満たすので$R$が正規でないことが分かる.\qed
\end{egans}











%-------------------局所環-----------
\subsection{局所環}

\begin{prop}
環$A$であって, $A$  がただひとつの極大イデアル$\maf{m}$を持つとき, 組$(A,\maf{m})$を局所環という. 局所環$(A,\maf{m})$について次が成り立つ.
\ben
\item \bolm{A^{\times} = A - \maf{m}}
\item イデアル$I\neq A$は$I\subset \maf{m}$を満たす,
\een
\end{prop}
\begin{egprob}
環$A$で, $A-A^{\times}$がイデアルであるとき, $A$は局所環であることを示せ.
\end{egprob}
\begin{egans}

\end{egans}
局所環は幾何的な文脈で自然に現れる. たとえば原点のある近傍で定義される$C^{\infty}$級関数の環$A$を考えよう. これは厳密に書くと$\varinjlim_{U\ni 0} C^{\infty}(U)$である. ただし, $U$は$0$を含む開集合を走る. すなわちこの環の元は, 開集合$U$と$C^{\infty}(U)$の元$f_{U}$の組$(U,f_U)$の 「0の十分近くで関数が一致する」という同値関係で割ったときの同値類$[U,f_U]$である. この環において, $[U,f_U]$がもし$f_U(0)\neq 0$であるなら任意の代表元に関して$f$の0での値は0でない. したがって$0$の近傍$V\subset U$で$1/f_U$が考えられ, $[V,(1/f_U)|_{V}]$が定まる. 明らかにこれは$[U,f_U] = [V,f|_{V}]$の逆元であるから, $A^{\times}$の元であることと$f(0)\neq 0$なる$f：U\to\mb{R}$で代表されることが同値であると分かる. 一方で$[U,f_U]$で$f_U(0)=0$を満たすものの全体を$\maf{m}\subset A$とおくと, 容易に分かるように$\maf{m}=A-A^{\times}$はイデアルをなす. したがって, 先の例題により$(A,\maf{m})$は局所環である. このように, \tf{関数環の茎}として現れるのが局所環である.

\begin{prac}
$(A,\maf{m})$をNoether局所環とする. $\bigcap_{n>0}\maf{m}^n = 0$を示せ.
\end{prac}


%-------------------Noether---------
\subsection{Noether環}
環の有限性条件である.

\tbox{
\begin{prop}
$A$を環とする。$A$が次の同値な条件のいずれか(したがってすべて)を満たすとき, $A$をNoether環という。
\ben[label=(\roman*)]
\item $A$のすべてのイデアルは有限生成である。
\item (\tf{asscending chain condition}) 任意の$A$のイデアルの昇鎖
\[I_1\subset I_2\subset \cdots \]
は停留する。つまり, ある番号$N$が存在して$I_N = I_{N+1} = \cdots$が成り立つ.
\item $A$の任意の空でないイデアル族が極大元を持つ.
\een
Noether環$A$に対して次が成り立つ。
\ben
\item (\tf{Hilbertの基底定理}) $A[x]$もNoether環である。
\item Noether環上の有限生成代数と局所化はNoether環である.
\item ($\spadesuit\spadesuit$, Krullの標高定理) $I$が$A$の$r$個の元で生成されるイデアルとするとき, $\mr{ht}(I) \leq r$である.
\een
\end{prop}
}{title=Noether環}
\prf (2): (1)が示せれば, 多項式環$A[x_1,\dots,x_n]$もNoetherで, その剰余環もNoetherである. 局所化がNoether環なのは, イデアルが拡大したものしかないことから明らか.\\
(1): これの証明は独特であり, しかし代数Iレベルなので覚えておこう. $\maf{a}\subset A[x]$が有限生成であることを示したい. $\maf{a}$から次の$A$のイデアル$I$を構成する:
\[I=\set{a\in A\mid aは\maf{a}の元の最高次係数として取れる}\]
これがイデアルであることはやさしい. よって$I$は有限生成で$I=(a_1,\dots, a_n)$と書ける. $f_i \in \maf{a}$が$f_i = a_ix^{r_i} + (\text{lower term})$と書けるとせよ. $r:=\max_{i=1}^{n}r_i$とおく. さて, $f\in \maf{a}$を取る. $f=ax^m + (\text{lower term})$とせよ. $a\in I$なので$a=\sum_{i} u_i a_i$と書ける. もし$m\geq r$であるなら,
\[f - \sum_{i} u_i f_ix^{m-r_i}\]
を考える. これは次数$m$未満の$\maf{a}$の元である. このように$f$から$\maf{b}=(f_1,\dots, f_n)\subset $の元を差し引いて, 次数$r$未満の元$g$を得ることができる. つまり,$f$に対してある次数$r$未満の$g$が存在し, $f-g\in \maf{b}$とできる. \par
$M=A1+Ax + \dots +Ax^{r-1}$とする. 先ほど示したことは$\maf{a} = (M\cap \maf{a}) + \maf{b}$である. $\maf{b}$は$A$上有限生成. $M\cap \maf{a}$も有限生成である. なぜなら, $M\cap \maf{a}\subset M$で, $M$が有限生成$A$加群で$A$はNoetherだからNoether $A$加群で, 部分加群$M\cap \maf{a}$も$A$加群として有限生成だから. \qed



\begin{prac}
Noether環の二つの同値な条件, の同値性を示せ.
\end{prac}

\begin{eg}
よく見る環はたいていNoetherである. たとえば$\mb{Z}, \mb{C}[X_1,\dots,X_n]$はNoetherであり, その剰余環や局所化もNoetherである.
\end{eg}

\begin{egprob} (仮想面接問題)
\ben
\item 非Noetherな環の例をあげてください.
\item 素イデアルが一つの環で非Noetherなものはありますか?
\een
\end{egprob}
\begin{egans}
(1): $k[x_1,\dots, ]$ \\
(2): $k[x_1,\dots]/(x_1^2,\dots)$
\end{egans}


\begin{rem}
Noether環の次元は有限とは限らない.  これは永田先生により反例が得られており, かなり難しい例がある.
\end{rem}

\begin{egprob}
Noether環の冪零元根基$\sqrt{(0)}$は冪零であることを示してください. もしNoether性を外すとどうなりますか?
\end{egprob}
\begin{egans}
$\sqrt{(0)}$が有限生成なので, その生成元が消える指数をすべて足したものでイデアルのべき乗を取ればよい. Noether性を外した場合, 次の例は$\sqrt{(0)}$が冪零ではない.
\[k[x_1,x_2,x_3,\dots],\quad (x_n^{n} = 0)\]
このことを見る. まず, この環の素イデアルは$(x_1,x_2,\dots)$のただひとつであることに注意. よって, このイデアル自身が冪零元根基である. ところが, このイデアルを$N$乗しても0にはならない. 実際, $x_{N+1}^{N}\neq 0$が残るから. \qed
\end{egans}


Noether局所環も重要である.
\begin{egprob}
Noether局所環$(A,\maf{m})$に対して, 次は同値であることを示せ.
\ben
\item $\dim{A} = 0$
\item $\maf{m} = \sqrt{0}$
\item $\maf{m}^q = 0$なる自然数$q$がある.
\item $A$はアルティン環.
\een
\end{egprob}
\begin{egans}
(i)$\implies$(ii): 冪零元根基が素イデアル共通部分なことと局所環であることから明らか.\\
(ii)$\implies$ (iii): $\maf{m}$が有限生成なことからただちに分かる. \\
(iv)$\implies$(i): アルティン整域は体だから. \\
(iii)$\implies$(iv): $k=A/\maf{m}$とする. $(I_n)_{n}$を$A$のイデアル減少列とする. $M_r = \maf{m}^r/\maf{m}^{r+1}$は$k$ベクトル空間である($r\geq q$では0である). このとき\bolm{(I_n\cap \maf{m}^r)/(I_n\cap \maf{m}^{r+1})}は$\maf{m}^{r}/\maf{m}^{r+1}$の$k$部分ベクトル空間の($n$に
関する)減少列をなす. $0\leq r\leq q$で$r$は有限通りであるから, 十分大きい$n_0$があり, $n\geq n_0$ならば$(I_n\cap \maf{m}^{r})/(I_n\cap \maf{m}^{r+1}) = (I_{n+1}\cap \maf{m}^{r})./(I_{n+1}\cap \maf{m}^{r+1})$である. $x\in I_n\cap \maf{m}^r$とするとき, ある$y\in I_{n+1}\cap \maf{m}^r$が存在して, $x-y \in I_n \cap \maf{m}^r$である. よって$I_{n}\cap \maf{m}^r \subset I_{n+1} +  (I_n\cap \maf{m}^r)$が従う. これがすべての$0\leq r\leq q$で成り立つので$I_n\cap \maf{m}^r\subset I_{n+1} + (I_n \cap \maf{m}^{q}) = I_{n+1}$となり$I_n - I_{n+1}$ ($n\geq n_0$)が得られた. \qed
\end{egans}

\begin{egprob}  \label{egprob:dimA_leq_m/m^2}
$(A,\maf{m})$をNoether局所環, $k=A/\maf{m}$とする. $\dim{A} \leq \dim_{k}\maf{m}/\maf{m}^2$を示せ.
\end{egprob}
\begin{egans}
$\maf{m}/\maf{m}^2 = kx_1 + \dots + kx_e$のとき$\maf{m} = Ax_1+\dots + Ax_e + \maf{m}^2$より, 中山の補題で$\maf{m}/ (Ax_1 + \dots + Ax_e) = 0$が言えるので$\maf{m}$は$e$元生成. よって標高定理から$\dim{A} = \mr{ht}(\maf{m}) \leq e$.\qed
\end{egans}





%---------------Artin----------------
\subsection{Artin環}

\tbox{
\begin{prop}
$A$を環とする。$A$が\tf{desscending chain condition}, すなわち「任意の$A$のイデアルの降鎖
\[I_1\supset I_2\supset \cdots \]
は停留する。つまり, ある番号$N$が存在して$I_N = I_{N+1} = \cdots$が成り立つ。」を満たすとき, $A$はArtin環という。Artin環について次が成り立つ。
\ben
\item Artin環の素イデアルは極大で, 素イデアルは有限個である。すなわち$\spec{A}$は離散位相空間である。
\item (\tf{秋月の定理}) Artin環であることと, 次元0のNoether環であることは同値である。
\item (\tf{Artin環の構造定理}) $A$はArtin局所環の積に同型であり, その積は一意的に定まる.
\een
\end{prop}
}{title= ($\spadesuit$) Artin環まとめ}


\begin{prac}
$\mb{Z}$はNoether環であることを示せ。$\mb{Z}$はArtin環ではないことをArtin環の定義から示せ.
\end{prac}
\begin{prac} (有名問題)
\tf{Artin整域は体であることを示せ.}
\end{prac}

\begin{rem}
$\spec{A}$が離散的でも$A$がArtin環であるとは限らない. ただし$A$がNoetherであるという亜k底をつければこれは同値である. 非Noetherだが素イデアルが一つしかないような例としては$\mb{C}[X_1,\dots, X_n,\dots]/(X_1^2,\dots, X_n^2,\dots,)$などがある.
\end{rem}

\begin{egprob}
Noether環$A$に関して, 次は同値であることを示せ.
\ben
\item $A$はArtin環
\item $\spec{A}$は離散位相空間
\een
\end{egprob}
\begin{egans}
(1)$\implies$(2)はよい. (2)$\implies$(1)を示す.  離散位相なのですべての点は閉点であり, すべての素イデアルが極大イデアルとなる. $A$はNoetherであるから, 秋月の定理より$A$はArtin環である.
\end{egans}


次はなかなか不思議な問題かもしれない.
\begin{egprob}
\ben
\item (\tf{有名}) $\spec{A}$がZariski位相でコンパクトであることを示せ.
\item $\spec{A}$が離散位相を持つなら, 有限集合であることを示せ.
\item 体$k$の可算直積環$\prod_{i=1}^{\infty} k$には$k\times k\times \dots k \times 0 \times k\times k\times \dots$以外の素イデアルが存在することを証明せよ.
\een
\end{egprob}




\begin{egprob}
$(A,\maf{m})$をNoether局所環とする. このとき, 次のいずれかが成り立つことを示せ.
\ben
\item すべての$k\geq 1$で$\maf{m}^k \neq \maf{m}^{k+1}$
\item ある$k\geq 1$で$\maf{m}^{k} = 0$.
\een
また, (2)が成り立つことと$A$がArtin局所環であることは同値であることを示せ. 秋月の定理は認めてよい.
\end{egprob}
\begin{egans}
(1)が成り立たないとせよ. もし$\maf{m}^k = \maf{m}^{k+1}$なら, 中山の補題($\maf{m}^k$は有限生成かつJacobson根基は$\maf{m}$)より$\maf{m}^{k} = 0$. このとき任意のイデアル$\maf{p}$は$\maf{p}\supset \maf{m}^k$で, $\maf{p}$が素だから$\maf{p}\supset \maf{m}$となる. よって$\maf{p} = \maf{m}$で$A$は0次元Noetherだから秋月の定理よりArtin. \par
逆に任意のArtin局所環を考える. この場合$\maf{m}$は冪零元根基でもあり, 有限生成であるから$\maf{m}^n = (0)$となる$n$が取れる(生成元が消える冪乗を総和すればよい).\qed
\end{egans}

%M1
次はできなくてもよいが, 証明が面白いので載せておく.
\begin{egprob}
Artin環がNoetherであることを用いずに, Artin環において冪零元根基$N=\mr{nil}A$は冪零であることを示せ.
\end{egprob}
\prf $N^k= N^{k+1} = \dots$なる$k$がある. $N^k \neq 0$と仮定する. $\maf{b}N^k\neq 0$を満たすすべてのイデアル$\maf{b}$の集合$\Sigma$は$N\in \Sigma$なので, Zornの補題から$\Sigma$の極小元$I$がある. $IN^k \neq 0$のとき, ある$x\in I$が存在して$xN^k \neq 0$なので(とくに$x\neq 0$かつ)極小性より$I=(x)$である. $xN^k \subset (x)$かつ$xN^k\in \Sigma$なので$xN^k=(x)$である. よって$x=xy$ ($y\in N^k$)と書ける. 一方, このとき$x=xy=xy^2=\dots = xy^n = \dots$で, $y\in N^k \subset N$だから$r$が十分大きいと$x = xy^{r} = 0$となり, $x$の選び方に反する.
\qed


\subsection{有限生成$k$代数}
体上の有限生成代数は, アファイン多様体の座標環に現れる環であり非常に重要である. とくに, Hilbertの零点定理が重要である.

\tbox{
\begin{prop}
\ben
\item (\tf{Noether正規化定理}) $A$を体$k$上の有限生成代数とする. つまり, $A=k[t_1,\cdots, t_r]$と書けるとする. このとき$T_1,\cdots, T_d\in A$という代数的独立\footnote{いかなる多項式関係式も満足しないという関係. このとき$k[T_1,\cdots, T_d]$は$d$変数多項式環に同型.} な元が存在して, $k[T_1,\cdots, T_d] \subset A$が整拡大であるようにできる. (ただし$d\geq 0$であり, $d=0$の場合は$k[T_1,\cdots, T_d] = k$と考える. )
\item (\tf{Zariski's Lemma}) $k$を体とする。有限生成$k$代数$A$の極大イデアルを$\maf{m}$とするとき, $A/\maf{m}$は$k$の有限次拡大体である。
\item (\tf{Hilbeltの弱零点定理}) とくに, $k$が\tf{代数閉体}ならば多項式環$k[T_1,\cdots,T_n]$の極大イデアルは$(T_1-\alpha_1,\cdots, T_n-\alpha_n)$という形をしている.
\item (\tf{Hilbertの強零点定理})$k$が\tf{代数閉体}であるとき, $\mac{I}\mac{Z}(I) = \sqrt{I}$
\een
\end{prop}
}{title=有限生成$k$代数まとめ}

どれも重要なのだが, \tf{正規化定理が一番強い.} それを例題として述べておく.
\begin{egprob}
$(1)\implies (2) \implies (3)$を示せ.
\end{egprob}
\begin{egans}
(1)を仮定せよ. $A/\maf{m}$は体であり, 有限生成$k$代数である. よって$k[T_1,\cdots, T_d]\subset A/\maf{m}$は整拡大だが, 右が体なので左も体である. よって$d=0$なので$A/\maf{m}$は$k$上整. つまり$k$上代数的. 有限次なのは有限生成だから. よって(2)が成り立つ. \\
(2)を仮定せよ. $k$を代数閉体とする. $k[T_1,\cdots, T_n]$の極大イデアル$\maf{m}$をとるとき, $k[T_1,\cdots, T_n]/\maf{m}$は$k$の有限次代数拡大である. しかし$k$は代数閉体なのでこれは$k$を超えることは出来ず, $k[T_1,\cdots, T_n]/\maf{m} = k$となる. よって, $T_i + \maf{m} = \alpha_i$なる$\alpha\in k$を見つけることができるので, $T_i - \alpha_i \in \maf{m}$である. よって$\maf{m}\supset (T_i - \alpha_i)_{1\leq i\leq n}$になる. 右は極大なので等号が成り立つ. よって示せた.\qed

\begin{egprob}
$A$が商体$K$を持つ整域で有限生成$k$代数であるとき, $\dim{A} = \mr{trdeg}_{k}K$であることを示せ. (スキーム論的に: $X$を整代数多様体とする. このとき任意の空でない部分代数多様体$U$に対して$\dim{X} = \mr{trdeg}_{k}K(X)$であることを示せ.)
\end{egprob}
\begin{egans}
\aka{正規化補題から単射整準同型$k[T_1,\dots,T_d] \hookrightarrow A$が存在.} $k[T_1,\dots,T_d]$の次元は$d$で, 整拡大の間で次元は不変なので(\ref{egprob:dimA_dimB_integral_ext}参照), $d=\dim{A}$である. この準同型を零イデアルで局所化して誘導されるものを考えると$k(T_1,\dots, T_n) \hookrightarrow K$を得る. これも体の間の単射整準同型であり, したがって代数拡大とみなせる. よって超越次数は$d$であるから示せた. \qed
\end{egans}




\begin{egprob}
$\sqrt{I} = \bigcap_{I\subset \maf{m}} \maf{m}$を示せ.
\end{egprob}
次はおぼえておくとよい.
\tbox{
\begin{egprob} \label{egprob:fg-k-alg-closedpt-to-closedpt}
$\phi:A\to B$を有限生成$k$代数の間の射とするとき, $\maf{n}\subset B$が極大ならば$\phi^{-1}(\maf{n})$も極大であることを示せ.
\end{egprob}
}{}
\begin{egans}
$\phi$は単射$A/\phi^{-1}(\maf{n})\to B/\maf{n}$を引き起こし, $B/\maf{n}$は$k$の有限次拡大. とくに$k$上の代数拡大であり, $k\subset A/\phi^{-1}(\maf{n})$とみなせるので$B/\maf{n}$は$A/\phi^{-1}(\maf{n})$上整である. よって$B/\maf{n}$が体だから$A/\phi^{-1}(\maf{n})$も体. \qed
\end{egans}


\end{egans}
次は専門題からの出題である.
\begin{egprob} (H29-専門2)
$K$を標数2でない代数閉体, $a\in K$とし
\[R_a = K[X,Y]/(X^2-Y^2-X-Y-a)\]
とおく.
\ben
\item $a=0$の場合に$R_a$の極大イデアル$\maf{m}$に対し$\dim_{K}\maf{m}/\maf{m}^2$を計算せよ. また, $\maf{m}$が単項イデアルになることはあるか.
\item $a\neq 0$の場合に上の問題を考えよ.
\een
\end{egprob}
\begin{screen}
$K$が代数閉体なので, $K=\mb{C}$の場合が基本である. $R_a$は曲線$X^2 - Y^2 - X - Y - a = 0$の上の関数環である. この曲線はよく見ると双曲線であるが, 中心がずれている. まずは, 平行移動するなどしてこの曲線を\tf{扱いやすい形に変える}ことが重要だ. 代数閉体なので$R_a$の極大イデアルは\tf{零点定理から曲線上の点に対応する.}  $\dim_{K}\maf{m}/\maf{m}^2$は, 代数幾何になじみがあれば答えが予測できるのだが, \tf{曲線の接空間の次元}になる.
\end{screen}
\begin{egans}
標数が2でないので$2^{-1}\in K$. よって, 次のような同型が作れる.
\bealn{
R_a &= k[X,Y]/(X^2-Y^2-X-Y-a) \\
&=k[X,Y]/((X-2^{-1})^2 - (Y+2^{-1})^2 - a) \\
&=k[T_1,T_2]/(T_1^2 - T_2^2 - a)\quad (T_1= X-2^{-1}, T_2 = Y+2^{-1})\\
&=k[T_1,T_2]/((T_1-T_2)(T_1+T_2) - a) \\
&\cong k[S_1,S_2]/(S_1S_2 - a)\quad (T_1+T_2 = S_1, T_1-T_2 = S_2)
}
$R_a$を最後の環と同一視する. $K$は代数閉体なので, 零点定理から$R_a$の極大イデアルは$xy = a$上の点$(x,y)\in K^2$と全単射対応する:
\[ (S_1-x, S_2-y)/(S_1S_2-a) \leftrightarrow (x,y)\in \set{xy=a}\subset K^2 \]
左を$\maf{m}$とおき, $S_{1,x} = S_1-x$, $S_{2,y} = S_2-y$とおく. $S_1S_2 - a = S_{1,x}S_{2,y} + xS_{2,y} + yS_{1,x}$である. $K$ベクトル空間として
\bealn{
\maf{m}/\maf{m}^2 &\cong (S_{1,x}, S_{2,y})/(S_{1,x}^2, S_{1,x}S_{2,y}, S_{2,y}^2, S_1S_2 - a)\\
&= (S_{1,x}, S_{2,y})/(S_{1,x}^2, S_{1,x}S_{2,y}, S_{2,y}^2, xS_{2,y}+yS_{1,x})\\
}
最後のベクトル空間で$f\in (S_{1,x},S_{2,y})$の類を$[f]$で表すと,
\bealn{
&= (S_{1,x}, S_{2,y})/(S_{1,x}^2, S_{1,x}S_{2,y}, S_{2,y}^2, xS_{2,y}+yS_{1,x})\\
&= K[S_{1,x}] + K[S_{2,y}] \neq 0
}
\\
となる. ここまでくれば, $[S_{1,x}]$と$[S_{2,y}]$が一次独立かどうかで$\dim_{K}\maf{m}/\maf{m}^2$が1か2かが決まる. 剰余加群の分母に$xS_{2,y} + yS_{1,x}$があるので, $x[S_{2,y}] + y[S_{1,x}] = 0$である. したがって, $(x,y)\neq (0,0)$なら$[S_{1,x}]$と$[S_{2,y}]$は一次従属であり, 逆に$(x,y)=(0,0)$なら一次独立である. $(x,y) = (0,0)$となるのは$a=0$のときだけなので, $\dim_{K}\maf{m}/\maf{m}^2$は以下の通り\footnote{曲線$xy=a$を$\mb{R}$上で図示すれば, 接空間の次元と一致することが確かめられる. }.
\[
\dim_{K}\maf{m}/\maf{m}^2 =
\begin{cases}
2\quad (a=0, x=y=0)\\
1\quad (otherwise)
\end{cases}
\]
%%%単項イデアル
$a=x=y=0$の場合, $\maf{m} = (s_1,s_2)$である($s_i = S_i \bmod{(S_1S_2)}$). もし$\maf{m} = (f(s_1,s_2))$となるとしよう. $f(s_1,s_2) = c + s_1f(s_1) + s_2f_2(s_2)$と一意に表せるが, 仮定より$c=0$が分かる. 条件より$(b + g_1(s_1)s_1 + g_2(s_2)s_2)(s_1f_1(s_1) + s_2f_2(s_2)) = s_1$となる$b\in k$, $g_1(s_1)\in k[s_1]$, $g_2(s_2)\in k[s_2]$が存在する. 右辺は$0+s_1\cdot 1 + s_2\cdot 0$だと思える. この形の表示は一意であることから考えれば, $b\neq 0$, $bf_1 = 1$ $f_1g_1 = 0$, $f_2g_2 = 0$が分かり, そこから$f_1\in k^{\times}$を得る. 同様に$s_2\in (f(s_1,s_2))$の表す式を書くことで$f_2\in k^{\times}$を得る. すると$f_ig_i = 0$に戻って$g_i = 0$を得るが, $bf_1 s_1 + bf_2s_2 = s_1$となるが, $bf_2 \neq 0$なのでおかしい. よって単項イデアルにならない. \par
それ以外の場合, $\maf{m} = (s_1 - x, s_2 - y)$である. $s_i := S_i \bmod{(S_1S_2 - a)}$である.\par
まず$a\neq 0$の場合は, 次の式から$\maf{m}$が単項イデアルになると分かる.
\[(-a)^{-1} xs_1(s_2 - y) =(-a)^{-1}(xa - as_1) = s_1 - x\]
最後に$a=0$の場合は, $x,y$のどちらか一方が0であり, $x\neq 0$なら次の式から$\maf{m} = (s_2)$が分かる.
\[s_2(s_1 - x) = - x s_2\]
\end{egans}

\begin{rem}
$\set{s_1s_2 = 0}$上の関数と見て考察するというのも考えられるが, これは標数0でない場合には危険である. $0$でない多項式がどこかの点で0でない値を取るとは限らないからだ.
\end{rem}

\begin{tips} より一般に$R$が代数曲線の座標環で$a$が曲線の点としよう. $R_{a}$の極大イデアル$\maf{m}$で局所化$A = (R_a)_{\maf{m}_{a}}$が考えられるが, このとき$\maf{m}A$が$A$の極大イデアルであり, $k=R_a/\maf{m}$ベクトル空間として$\maf{m}A / \maf{m}^2A \cong \maf{m}/\maf{m}^2$である. このときの$\maf{m}A/\maf{m}^2A$の$k$上の双対空間をZariski接空間といい, 一般に\tf{曲線の非特異点では$\dim{A} = \dim{\maf{m}/\maf{m}^2}$}が成り立つ.  このときの$\maf{m}A$が\tf{何項で生成されるか}という情報は実は重要で, 非特異点$a$の局所環$A$では$\maf{m}A$が$d=\dim{A}$項生成になるというのがある. いまは曲線だから非特異点で$d=1$であり, $\maf{m}A$は単項生成になる. $\maf{m}$と$\maf{m}A$はもしかしたら違うかもしれないが\footnote{正確な情報求む.}, こういう背景であたりをつけることができる. 逆に特異点では$\dim{A} < \dim{\maf{m}/\maf{m}^2}$となり, $\maf{m}$が$d=\dim{A}$項より多くないと生成できないことがある.
\end{tips}

\tbox{
\begin{prac} (仮想面接問題)\\
$A$を$d$次元正則局所環であるとき, $\maf{m}$は$d$個の元で生成できることを示せ.  (\tf{Hint:中山の補題})
\end{prac}
}{}
% 正則局所環なので$\dim_{k}\maf{m}/\maf{m}^2 = d = \dim{A}$である. よって$\maf{m}/\maf{m}^2 = kx_1+\dots + kx_d$となる$x_1,\dots, x_d\in \maf{m}/\maf{m}^2$がある. $x_i$が$y_i \bmod{\maf{m}^2}$ ($y_i \in \maf{m}$)であるとすると, $\maf{m} = Ay_1+\dots + Ay_d + \maf{m}^2$が集合論的に示せる. このとき, $M = Ay_1+\dots + Ay_d$として$\maf{m}/M = \maf{m}(\maf{m}/M)$であり$A$は局所環だから中山の補題より$\maf{m} = M$である. \qed



















これから述べる環は発展的である.



\subsection{Dedekind環}

\begin{prop} (\tf{Dedekind環まとめ})\\
1次元のNoether整閉整域を\tf{Dedekind環}という。Dedekind環$A$について次が成り立つ。
\ben
\item 素イデアル分解が可能である.
\item 任意の素イデアル$\maf{p}$に対して$A_{\maf{p}}$は離散付値環である.
\een
\end{prop}
こと代数的整数論では整数環を扱う.
\begin{theo}
代数体$K$の整数環$\mac{O}_K$はDedekind環である.
\end{theo}

\begin{recall} $d\neq 0$が平方因子を持たない整数とする.
$\mb{Q}(\sqrt{d})$の整数環は$d\equiv 1\pmod{4}$で$\mb{Z}[(\sqrt{d}+1)/2]$で$d\equiv 2,3\pmod{4}$なら$\mb{Z}[\sqrt{d}]$.
\end{recall}

素イデアル分解が口頭試問で問われる可能性は高いだろう. 「上にある素イデアル」を求めるのが筋である.



\begin{egprob}
$\mb{Z}[\sqrt{-5}]$において, $(6)$の素イデアル分解は?
\end{egprob}

\begin{egans}
$(2),(3)$を分解して得られる. $\mb{Z}[\sqrt{-5}]/(2)\cong \mb{F}_{2}[X]/((X-1)^2)$より$\maf{p}_{2}=(2,1+\sqrt{-5})$は$\maf{p}_{2}^2 = (2)$. $\mb{Z}[\sqrt{-5}]/(3)\cong \mb{F}_{3}[X]/(X^2-1)$より$\maf{p}_{3,1} = (3,\sqrt{-5} + 1)$, $\maf{p}_{3,2} = (3,\sqrt{-5}-1)$が$\maf{p}_{3,1}\maf{p}_{3,2} = (3)$. よって$(6) = \maf{p}_{2}^2\maf{p}_{3,1}\maf{p}_{3,2}$.
\end{egans}

\begin{recall}
\tf{$\mb{Z}[\sqrt{-5}]$はUFDでない整閉整域の例. }
\end{recall}

\subsection{付値環} %M1に追加
院試では見たことがないが, \cite{atimac}の5章にある.
\begin{defn}
整域$B$と商体$K$が, 「$x\in K$なら$x\in B$か$x^{-1}\in B$のどちらか少なくとも一方が成立する」を満たすとき, $B$を\tf{付値環}という.
\end{defn}

\begin{prop} 付値環$B$に対し,次が成り立つ.
\ben
\item $B$は局所整閉整域である.
\een
\end{prop}
\prf
(1): \\
\fbox{局所環であること}$\maf{m} = B-B^{\times}$がイデアルをなすことを示せばよい(単元を含まない最大の集合なので明らかに唯一の極大イデアルである). $a\in B, x\in \maf{m}\implies ax\in \mac{m}$は簡単. $x,y\in \maf{m}$のとき, $x+y = x(1+x^{-1}y) = y(xy^{-1} + 1)$に注意し, $xy^{-1}$か$x^{-1}y$が$B$に属すことに注意すれば, $x+y\in\maf{m}$がわかる. \\
\fbox{整閉整域であること} \\
$x\in K$を$B$上整とする. $x\notin B$だとせよ. このとき$x^{-1}\in B$である. 整等式
\[x^n + b_1x^{n-1} + \dots + b_n = 0\]
に$x^{1-n}$をかけると$x\in B$が分かるがこれは矛盾. \\
(2):




\subsection{正則局所環}
代数幾何で正則性を定義するにあたって重要な環である.
\tbox{
\begin{defn}
$A$が正則局所環であるとは, Noether局所環$(A,\maf{m})$であって, $k=A/\maf{m}$とするとき, Noether局所環では一般に成り立つ「$\dim_{k} \maf{m}/\maf{m}^2 \geq \dim{A}$」の等号が成り立つような環のことをいう. $d$次元正則局所環$(A,\maf{m})$について, 次が成り立つ.
\ben
\item \tf{$\maf{m}$は$d$個の元で生成される. } また, これがNoether局所環に対する, $d$次元正則局所環であることと同値な条件である.
\item $A$は整域である.
\item $A$はUFDである. (\tf{これはかなり難しい命題})
\item $\maf{p}\in\spec{A}$に対し$A_{\maf{p}}$も正則.
\een
\end{defn}
}{title=$(\spadesuit\spadesuit)$ 正則局所環}






Noether環の次元論については様々なことが知られている. 証明はしない.
\begin{prop} (\cite{liu}2.5.2節などを参照)
$A$をNoether環とする.
\ben
\item $\maf{m}$が極大イデアルで$\maf{n}$が$A[T_1,\dots, T_n]$の極大イデアルで$\maf{n}\cap A = \maf{m}$を満たすものとする. このとき\bolm{\mr{ht}(\maf{n}) = \mr{ht}(\maf{m}) + n}
\item \bolm{\dim{A[T_1,\dots, T_n] = \dim{A} + n}}
\een
\end{prop}

\begin{eg}
\ben
\item $\maf{n}$を$\mb{Z}[T_1,\dots, T_n]$の極大イデアルとする. 実は$\maf{n}$はある素数を含むので$\mr{ht}(\maf{n}) = n+1$.
\item $\maf{m}$を$k[T_1,\dots, T_n]$ ($k$は体) の極大イデアルとする. $\maf{m}\cap k = (0)$なので$\mr{ht}(\maf{m}) = n$であるから$\dim{A_{\maf{m}}} = n$である.  また, $\maf{m}$は$n$元で生成されることが知られているので$\maf{m}A_{\maf{m}}$は$n$個の元で生成される. よって$\maf{m}A_{\maf{m}}/\maf{m}^2A_{\maf{m}}$ は$A_{\maf{m}}/\maf{m}A_{\maf{m}} = k$ベクトル空間として$n = \dim{A_{\maf{m}}}$個の元で生成され, (\ref{egprob:dimA_leq_m/m^2})により基底にもなっていなければならない. よって$\dim{A_{\maf{m}}} = \dim_{k}\maf{m}A_{\maf{m}}/\maf{m}^2A_{\maf{m}}$が成り立つので, $A_{\maf{m}}$は\tf{正則局所環}である.
\een
\end{eg}





\subsection{離散付値環}
代数曲線の非特異点の局所環として現れるのが離散付値環であり, その極大イデアルは一意化元の役割を果たす. 雰囲気としてはかなり, リーマン面の関数環の局所環のように思える. 1次元正則局所環なので, 正則局所環の性質も持っているが, 離散付値環はその次元の小ささから比較的扱いやすい.

\tbox{
\begin{defn}
1次元のNoether整閉整域で局所環であるもの, あるいは「離散付値」の付置環であるものを離散付値環(DVR)という. DVR$(A,\maf{m})$に対し, 次が成り立つ.
\ben
\item
\item \tf{1次元正則局所環と離散付値環は同じである. }とくに, $\maf{m}$が1つの元$t$で生成される.
\item $A$の任意の0でないイデアルは$(t^n)$ ($n\geq 0$) の形に表せる. とくに, \tf{DVRはPIDである.}
\een
\end{defn}
}{title= ($\spadesuit\spadesuit$)DVRまとめ}

\begin{eg}
$\mb{Z}_{(p)}$は$p\mb{Z}_{(p)}$を極大イデアルとする離散付値環である. 実際, 1次元であることはよい. Noether, 整閉も局所化で不変だからよい. 任意のイデアルは$p^n\mb{Z}_{(p)}$の形で表せる. 要は, 分子が$p$で何回割り切れるかでイデアルが分類できる.
\end{eg}


\begin{egprob} (H12数学II-2) \label{prob:H12-II-2}
離散付値環$(R,\maf{m})$とその一意化元$t\in \maf{m}$を考える. (この問での離散付値環の定義は1次元局所Noether整閉整域である) $K=\mr{Frac}R$とする.
\ben
\item 自然数$n$に対し$X^n-t$は$K[X]$で既約であることを示せ.
\item $R[X]/(X^n-t)R[X]$は離散付値環になることを示せ.
\een
\end{egprob}
\begin{egans}
(1): DVRはUFDで, $t$は素元で, $X^n - t$はアイゼンシュタイン多項式なので既約. \\
\fbox{ちゃんと書きたい場合} $X^n - t$が$K$上可約であるとする. $X^n - t = f(X)g(X)$とし, $f,g\in K[x]$はモニックかつ次数は1以上とする. $R$はUFDなので, 最小公倍元が定義できる. $f,g$の係数の分母の最小公倍元をそれぞれ$a,b$とすると, $af(X)$と$bg(X)
$は$R$の原始多項式であり, $ab(X^n - t) = (af(X))(bg(X))$である. 右辺は\tf{原始多項式の積だからまた原始多項式}であり, $ab$は素元を持ち得ない. よって$ab\in R^{\times}$だから, $a,b\in R^{\times}$であり, もともと$f,g\in R[X]$であることが分かる (よって$X^n-t$は$R[X]$でも可約). このとき$X^n \equiv f(X)g(X) \bmod{(t)}$である. $f(X) \bmod{t},g(X)\bmod{t} \in (R/tR)[X]$の最低次の部分に注目すれば, $f(X),g(X)$はそれぞれ, 最高次以外の項が$t$で割れることが分かる. $f,g$は一次以上だから, とくに定数項がそれぞれ$t$で割れるが, これは$X^n-t$の定数項が$t$であることに反する. よって既約である. \qed \\
(2): $A=R[X]/(X^n-t)$とする. $A$がNoetherは自明. $R[X]$はUFDであるから, $R[X]$の既約元は素元でもあり, したがって$A$は整域である. \\
\fbox{整閉性} $L=\mr{Frac}A = K(u)$ ($u^n = t$)である. $u$は$R$上整($K$上代数的)なので$L=K[u]$である. (1)より$[L:K] = n$である. $L$の元は$\alpha = \sum_{j=0}^{n-1} a_jt^{n_j} u^j$と書ける. ただし$a_j\in R, n_j\in \mb{Z}$.  いまこの元が$A=R[u]$上整であるとする(予想: $\alpha \in A$). $n_j$の最小値を$m$とし, $n_{i_1},\dots, n_{i_r} = m$であるとする. \ao{仮に$m<0$であったとせよ.}
$\alpha$に$t^{-m-1}\in R$をかければ, $n_j > m$であるような$u^j$の係数が$R$の元になり, $n_j = m$であるような$u^j$の係数は$\dfrac{a_j}{t}$になる. $\alpha t^{-m-1}$もまた$R[u]$上整なので, $R[u]$の元になった部分を差し引くことで, $\beta:=\sum_{j=1}^{r} \dfrac{a_{i_j}}{t}u^j\in L$が$R[u]$上整である. $i_1<\dots < i_r$であるようにする. \ao{$\beta$に$u^{n-1-i_1}$回掛けることで, $\dfrac{a_{i_1}}{t}u^{n-1}$が$R[u]$上整であることが分かる. }これを$n$乗して$\gamma:=\dfrac{a_{i_1}^n}{t^n}u^{n(n-1)} = \dfrac{a_{i_1}^n}{t}$が$R[u]$上整である. $R\subset R[u]$は整拡大だから$\gamma\in K$は$R$上整である. しかし$\gamma$は明らかに$R$に属しておらず, $R$が整閉であることに反する. したがって$m\geq 0$なので$A$が整閉であることが示された. \\
\fbox{局所環であること} $R\subset A = R[u]$は整拡大であるから, $f:\spec{A}\to \spec{R}$は閉点を閉点に移す($\because R/f(\maf{p})\subset A/\maf{p}$は整拡大なので, $R/f(\maf{p})$が体$\iff$ $A/\maf{p}$が体). いま$A$の極大イデアル$\maf{m}$を任意に取る. $f(\maf{m}) = \maf{m}\cap R$は$\spec{R}$の閉点だから$f(\maf{m}) = tR$である.  よって$\maf{m}$は一意化元$t = u^n$を含むから, $u\in \maf{m}$である. よって $\maf{m} \supset uA$が従う. $uA$は実は極大イデアルである. 実際,
\[A/uA \cong \dfrac{R[X]}{(X^n-t)} / \dfrac{(X,X^n-t)}{(X^n-t)} \cong R[X]/(X,t) \cong R/tR\]
は体であるからだ. よって$\maf{m} = uA$となり, $A$が局所環であることが示せた. \\
\fbox{1次元であること} $R\subset A$が整拡大であることを用いれば, $\dim{A} = \dim{R} = 1$なのでよい.\\
\fbox{1次元であること(ちゃんとやる場合)} 例の「整拡大ではファイバーの点同士に埋没がない」の証明を真似る. $\dim{A} \geq 1$はよいだろう. $R\subset A$は整拡大である. $\maf{p}\in\spec{A}$という極大でない素イデアルを任意に取る. $R\subset A$が整拡大であることから, $f(\maf{p})$は閉点にはならない. よって$f(\maf{p}) = 0R$である. ここで$R\subset A$が整拡大なので$R_{0R} = K\subset A_{\maf{p}}$も整拡大. $K$は体だから$A_{\maf{p}}$も体だが, このとき$\maf{p}A_{\maf{p}} = 0 = 0A_{\maf{p}}$である. よって局所化における素イデアル対応定理から$\maf{p} = 0A$を得る. したがって次元は1である. \qed

\end{egans}


\subsection{Cohen-Macaulay環}
可換環論で最も重要と言われているが, アティマクにも載っていない. 某IUTの教授が\kenten{アティマクからの質問として}CM環について聞いてきたという三次情報がある.

\begin{defn}
Cohen-Macauley環
\end{defn}







\subsection{カテナリー環}
素イデアルの減少列
\[\maf{p}_0\supsetneq \maf{p}_1\supsetneq \dots \supsetneq \maf{p}_n\]
\[\maf{q}_{0} \supsetneq \maf{q}_1\supsetneq  \dots \supsetneq \maf{q}_m\]
に関して, $\maf{p}_0 - \maf{q}_0$, $\maf{p}_n = \maf{q}_m$かつ$\maf{p}_i$が$\maf{q}_j$のどれかに等しいとき, 下の列は上の列の\tf{細分}であるという.

\begin{prop} $A$の任意の素イデアル$\maf{p}\subset \maf{q}$に対し, $\maf{p}$から始まり$\maf{q}$に終わる素イデアル減少列であって細分のできないようなものの長さが同じ(ただしその長さは$\maf{p},\maf{q}$に依存)であるとき,
\end{prop}






\subsection{演習問題}

\subsubsection{Level A}

\subsubsection{Level B}
\begin{prob}
Noether局所環$(A,\maf{m})$について, $\bigcap_{n\geq 0} \maf{m}^n = 0$を示せ.
\end{prob}


\begin{prob} (H1数学II-1)
単位的可換環$R$に関する次の命題1,2,3のおのおのについて, 次の問に答えよ.
\ben
\item 命題が正しいかどうか, 理由をつけて答えよ.
\item 命題が正しくない場合には, $R$に適当な条件をつけ加えて正しい命題にし, それを証明せよ.
\een
\begin{enumerate}[label=命題\arabic*]
\item $a_1,\dots, a_n\in R$, $f(x) = x^n + a_1x^{n-1} + \dots + a_n$に対し, $\set{c\in R\mid f(c) = 0}$の元数は$n$以内である.
\item $a\in R$が$R$の拡大環の中に逆元$a^{-1}$を持っていて, さらに
$R[a^{-1}]$が$R$加群として有限生成であれば$a^{-1}\in R$.
\item $R$上の一変数多項式環$R[x]$において逆元をもつ元は$R$の元に限る.
\end{enumerate}
\end{prob}



\subsubsection{Level C}


\subsubsection{Hint・Answer}


\begin{probans}
Artin-Reesの補題から中山の補題. なお, $\maf{m}$を真のイデアルにしても成り立つ.
\end{probans}


\begin{probans}
(1): \fbox{命題1}整域ではない$R$に対してこれは成り立たない. $ab=0$, $a,b\neq 0$として, $f(x) = x(x+b-a)$とする. これは$x=0,a,a-b$を根に持つが, $0\neq a-b$でないならこれは3種類になる. そのような整域はあるから, 正しくない. \\
\fbox{命題2} $a^{-1}$が$R$上整である. 整等式を立てれば良い. 正しい. \\
\fbox{命題3} 偽. $R$が被約でないとき, $a^N = 0$なる$a$に対し$aX - 1$は可逆(等比級数を考えよ).\\
(2): \fbox{命題1}$R$が整域のときに正しい. 次数に関する帰納法と除法を考えれば良い. \\
\fbox{命題3} アティマクの演習を復習しよう. $f(x)$が可逆であるための必要十分条件は, $a_0$が可逆でかつ$a_i$ ($i\geq 1$)が冪零であること. 単元と冪零元の和は単元であることに注意.
\end{probans}





\newpage

\section{可換環の理論の復習}
以降$A$は環とする。
\begin{prop} (対応定理)\\
$A$を可換環とする。$A$の$I$を含むイデアル$J$と, $A/I$のイデアルは1:1に対応する。
\end{prop}
\begin{prop}
$A$の任意のイデアル$I\subsetneq A$に対して, $I$を含む極大イデアルが存在する。
\end{prop}



\begin{theo} (\tf{中国剰余定理}, CRT)

\end{theo}
\begin{theo} (\tf{極大イデアルの存在})

\end{theo}

\begin{theo} (\tf{Hilbertの基底定理})

\end{theo}

\begin{theo}

\end{theo}



代数系の人が次の(i)を聞かれたという噂がある.
\begin{egprob} (\cite{atimac} Ex1.2) \\
$f=a_0+a_1x+\dots +a_nx^n\in A[x]$として, 次を示せ.
\ben
\item $f$が$A[x]$の単元である$\iff$ $a_0\in A^{\times}$で$a_1,\dots, a_n \in \sqrt{0A}$.
\item $f\in \sqrt{0A[x]} \iff a_0,a_1,\dots, a_n \in \sqrt{0A}$.
\item $f$が零因子$\iff$ $\exists a\in A\setminus{\set{0}}$, $af = 0$
\item $(a_0,a_1,\dots, a_n) = (1)$であるとき, $f$は原始的であるという. $f,g\in A[x]$のとき, $fg$が原始的$\iff$ $f$と$g$が原始的.
\een
\end{egprob}
\begin{egans}
(1). $g=b_0 + b_1x + \dots + b_mx^m$が逆元とせよ. $b_ma_n = 0$, $b_{m}a_{n-1} + b_{m-1}a_n = 0$より$a_{n}^2 b_{m-1} = 0$を得る. 次に$b_{m}a_{n-2} + b_{m-1}a_{n-1} + b_{m-2}a_{n} = 0$なら$a_{n}^3 b_{m-2} = 0$を得る.  これを繰り返し, $a_{n}^{m+1}b_{0} = 0$を得る. $b_0a_0 = 1$で$b_0$は単元なので$a_{n}\in \sqrt{0A}$が示せた. よって$a_{n}x^n \in \sqrt{0A[x]}$ なので, \aka{$f-a_{n}x^n$は単元と冪零元の和であり, これは単元である.} よって, $n-1$次以下の主張に帰着されるから,帰納的に示すことができる. なお, $0$次の場合の主張は明らかである. $(\Leftarrow)$は, まさしく$a_1x+\dots + a_nx^n$が冪零元で$a_0$が単元であることから単元になる. \\
(2) $\Leftarrow$はやさしい. 逆を示す. $a_n\in \sqrt{0A}$がすぐ分かるので, $a_{n}x^n\in \sqrt{0A[x]}$である. よって$\sqrt{0A}$はイデアルだから$f-a_nx^n\in \sqrt{0A[x]}$. これを続ければ$a_{0},\dots, a_{n-1}\in \sqrt{0A}$であることも言える.
\end{egans}


7
\subsection{($\spadesuit$)発展的内容をいくつか}
\subsubsection{(Local Property}
可換環論および加群論の多くの場面で, $A$のすべての$\maf{p}$での局所化$A_{\maf{p}}$がある性質$\hana{P}$を持つことと$A$が性質$\hana{P}$を持つことが同値である。このような性質$\hana{P}$を局所的性質という。局所化もそうだが, この局所的という意味は, $A$がアファインスキーム$\spec{A}$上の正則函数環と思えるのに対し, $A_{\maf{p}}$は$\maf{p}\in \spec{A}$での十分近くの正則函数環と思えることから来ている。
\begin{prop} (\tf{環のLocal Propertyまとめ})
以下はLocal Propertyである。
\ben
\item 整閉性. つまり, $A$が整閉$\iff$ $\forall \maf{p}\in \spec{A}$, $A_{\maf{p}}$が整閉
\item
\een
\end{prop}
\begin{eg}
「体である」という性質はLocal Propertyかどうか調べよう。$\forall \maf{p}, A_{\maf{p}}$が体とする。$A\neq 0$なので$\maf{p}$が少なくとも一つ存在する。このとき$\maf{p}A_{\maf{p}}$は$\maf{p}\in \spec{A}$に対応するが, $A_{\maf{p}}$が体なので$\maf{p}^{e} = \maf{p}A_{\maf{p}} = (0)$. $\maf{p} = \maf{p}^{ec} = \maf{p}A_{\maf{p}}\cap A = (0)$なので$\maf{p}=(0)$である。よって$A$の素イデアルは$(0)$のみである。$x\in A\setminus{\{ 0 \}}$かつ$x$が可逆でないならば, $x$を含む極大イデアルが存在するはずなので矛盾。よって$x$は可逆であり, $A$は体である。逆に体$A$の素イデアルは$(0)$しかなく, $A$と$A$の局所化は同じなのですべての素イデアルに対して局所化は体である。したがってLocal Propertyである。
\end{eg}
\begin{prac}
$A=\mb{C}[X,Y]/(X^2,XY,Y^2)$の素イデアルを決定し, この環を考えることで, 被約性はLocal Propertyではないことを示せ。ただし環$A$が被約であるとは, $\mr{nil}{A} = \bigcap_{\maf{p}\in \spec{A}}\maf{p}$が零であることをいう。
\end{prac}
%Aの素イデアルは(X),(Y),(X,Y)で被約だが(X)での局所環A_{(X)}は素イデアルが(X)のみなのでnil(A_{(X)})は(X)A_{(X)}であり, これは零ではない。たとえばX/(X-a)  (aは0でない複素数)はベキ零元。

\begin{egprob} (\tf{有名問題！}) \\
整域であることはLocal Propertyであるか?
\end{egprob}
\begin{egans}
\tf{No!} たとえば$A=\mb{Q}\times \mb{Q}$を考える. $\maf{p} = \mb{Q}\times 0$と$\maf{q} = 0\times \mb{Q}$がこの環の素イデアルのすべてであるが, このとき$A_{\maf{p}} \cong \mb{Q}$なので素イデアル局所化は整域である. しかし$A$は整域ではない. \\
\underbar{$A_{\maf{p}}\cong \mb{Q}$について}

\end{egans}













\subsection{例題集}

\begin{egprob}
環$R_1,R_2$に対して$R=R_1\times R_2$を直積環とするとき, 集合として
\[R^{\times} = R_1^{\times} \times R_2^{\times}\]
を示せ.
\end{egprob}
\begin{egans}
$R^{\times}$の元は, とりあえず$(r_1,r_2)$と書ける($r_i \in R_i$). これが可逆だから$(r_1,r_2)(s_1,s_2)=(1,1)$となる$s_i\in R_i$が存在する. 左辺は$(r_1s_1, r_2,s_2)$だから$s_i = r_i^{-1}$であり$r_i \in R_i^{\times}$である. 逆に$r_i\in R_i^{-1}$を与えれば, $(r_1,r_2)$は$R$で逆元$(r_!^{-1}, r_2^{-1})$を持つので$(r_1,r_2)\in R^{\times}$である. \qed
\end{egans}
\begin{eg}
$n\in \mb{N}$を素因数分解して$n=\prod_{j=1}^{d} p_j^{e_j}$とするなら, 環の中国剰余定理により
\[\mb{Z}/n\mb{Z}\cong \prod_{j=1}^{d} \mb{Z}/p_j^{e_j}\mb{Z} \]
なので
\[(\mb{Z}/n\mb{Z})^{\times} \cong \parena{\prod_{j=1}^{d}(\mb{Z}/p_j^{e_j} \mb{Z})}^{\times} = \prod_{j=1}^{d} (\mb{Z}/p_j^{e_j}\mb{Z})^{\times}\]
$\mb{Z}/N\mb{Z}$の位数はトーシェント関数により$\phi(N)$だから, 上の位数を比較することで, $\phi$の乗法性
\[\phi(n) =  \prod_{j=1}^{d} \phi(p_j^{e_j}) = \prod_{j=1}^{d} p_j^{e_j} \parena{1-\dfrac{1}{p_j}} = n\disp\prod_{j=1}^{d} \parena{1-\dfrac{1}{p_j}}\]
が導かれる. たとえば$\phi(48) = 48(1- \frac{1}{2})(1-\frac{1}{3}) = 16$である.
\end{eg}

\begin{egprob} (\cite{dummit}, 258p, 25)\\
$R$を可換環とする。任意の$a\in R$に対して$n_{a}$という$a$に依存する2以上の自然数が存在するとき, $R$の素イデアルはすべて極大イデアルであることを示せ。
\end{egprob}
\begin{egans}
$\maf{p}\subset R$を素イデアルとする。$\maf{p}$が極大でないとすると, $\maf{p}$を真に含む極大イデアル$\maf{m}$が存在する。$a\in \maf{m}\setminus{\maf{p}}$とする。$a^{n_{a}} = a$だから$a(a^{n_{a} - 1} - 1)=0\in \maf{p}$である。$a\notin \maf{p}$と$\maf{p}$が素イデアルでないことから$a^{n_{a}-1}-1\in \maf{p}\subset \maf{m}$である。ところが$a^{n_{a}-1}\in \maf{m}$であるから$1\in \maf{m}$となり矛盾である。\qed
\end{egans}

\begin{egprob} \label{egprob:Q_is_not_finitelygenerated_over_Z}
$\mb{Q}$は有限生成$\mb{Z}$代数ではないことを示せ.
\end{egprob}
\begin{egans}
混同しやすいことだが, 有限生成$A$加群と言われたら$\sum Ax_n$と書ける加群のことをいう. 有限生成$A$代数は, $A$係数有限変数多項式環の準同型像のことである. つまり(準同型定理からすると) イデアル$I\subset A[T_1,\cdots, T_n]$によって$A[T_1,\cdots,T_n]/I$で表せる環のことをいう. 有限生成$A$代数には, 当然多項式環$A[T]$にも含まれるわけだから, 有限生成$A$加群であるとは限らない.  \par
本題に入る. もし$\mb{Q} = \mb{Z}[p_1,\cdots, p_N]$ ($p_i\in \mb{Q}$) と書けるなら, 既約分数表示における$p_1,\cdots, p_N$の分母に表れない素数$q$をひとつ選べば, $\frac{1}{q} \notin \mb{Z}[p_1,\cdots, p_n]$が示せるのでおかしい. \qed
\end{egans}

\begin{egprob} (有名問題)
位数有限の整域は体であることを示せ。
\end{egprob}
\begin{egans}
$A$を位数有限の整域とする。有限性から$a\in A\setminus{\{ 0 \}}$に対して$a^{n} = a^{m}$を満たす自然数$n,m$ ($n>m$)が存在する。$a^{m}(a^{n-m} - 1) = 0$だから整域の性質から$a^{n-m} = 1$または$a^{m} = 0$である。$a\neq 0$と整域の性質から後者は生じないので$a^{n-m} = 1$である。よって$a^{n-m-1}$が$a$の逆元であるから$A$は体である。
\end{egans}


\begin{egprob}
次の非同型を示せ。
\ben
\item $\mb{Z}\not \cong \mb{Z}[\sqrt{-1}]$
\item $\mb{C}\not \cong \mb{R}$
\item $\mb{C}[x,y]/(xy)\not\cong \mb{C}\times \mb{C}$
\item $\mb{C}[x,y]/(xy)\not\cong \mb{C}[x]\times\mb{C}[x]$
\een
\end{egprob}


\begin{egprob}
$\mb{Z}_{(2)}[T]$において, 極大イデアルは必ず2を含むか?
\end{egprob}
\begin{egans}
反例: $(2T-1)$. このイデアルで割ると$\mb{Q}$になるので確かに極大である.
\end{egans}
\begin{egans}
Jacobson根基(すべての極大イデアルの共通部分)を使って答えることもできる. もし$2$がそれに入るなら, $1+2T$は$\mb{Z}_{(2)}[T]$の可逆元である(Why?). これはおかしい.
\end{egans}





\begin{egprob} (9/12, Quiz, \cite{aoyukie} Ex1.4.1)
\ben
\item $\mb{Z}[\sqrt{2}]/(3+\sqrt{2}) \cong \mb{F}_{7}$であることを示せ。
\item $\mb{Z}[\sqrt{5}]/(1+2\sqrt{5})$ は何か? (簡単な環と同定せよ)。
\een
\end{egprob}
\begin{egans}
(1): 同型$\mb{Z}[x]/(x^2-2) \cong \mb{Z}[\sqrt{2}]$でイデアル$(3+\sqrt{2})$はイデアル$(3+x,x^2 - 2)/(x^2 - 2$に対応するので
\beqn
&& (\mb{Z}[x]/(x^2 - 2))/((3+x,x^2-2)/(x^2-2))\\
&\cong& \mb{Z}[x]/(x+3, x^2-2) \\
&\cong & \mb{Z}/7\mb{Z}
\eeqn
ただし, 最後の同型は$[x] \mapsto -3 \bmod{7}$による($x=-3$の代入$\mb{Z}[x] \to \mb{Z}$から誘導される準同型)。

(2) $\mb{Z}[\sqrt{5}] \cong \mb{Z}[x]/(x^2 - 5)$と$4x^2 - 20 = -(1+2x)(1-2x)-19$より
\beqn
&&\mb{Z}[\sqrt{5}]/(1+2\sqrt{5}) \\
&\cong& (\mb{Z}[x]/(x^2 - 5)) / (1+2x, x^2-5)/(x^2 - 5)\\
&=& (\mb{Z}[x]/(x^2-5)) / ( (1+2x,\textcolor{blue}{4x^2 - 20}, x^2 - 5)/(x^2-5))\\
&\cong& \mb{Z}[x] / (1+2x, -19, x^2 - 5)\\
&=& \mb{Z}[x]/(1+2x, 19, x^2-5) \\
&\cong& (\mb{Z}[x]/(19))/(1+2x, x^2-5)\\
&\cong& \mb{F}_{19}[x] / (1+2x, x^2-5)\\
&=& \mb{F}_{19}[x]/(10+x,x^2-5)\quad (\because 10(1+2x) = x+10)\\
&\cong& (\mb{F}_{19}[x]/(x+10))/(10+x,x^2-5)/(10+x)\\
&\cong& \mb{F}_{19}/(19, 9^2 - 5)\qquad (\because x=9の代入)\\
&=& \mb{F}_{19}/(0) \\
&=& \mb{F}_{19}
\eeqn
\end{egans}


\tbox{
\begin{egprob}
\ben
\item PIDではないUFDの例は?。
\item Noetherでない環の例は?
\item ArtinでないNoether環は?
\item UFDでない整閉整域の例は?
\een
\end{egprob}
}{title=環のクラスと反例}
\begin{egans}
(1):$\mb{Z}[x], \mb{R}[x,y]$など。まずUFDの多項式環はUFDであるから, この両者はUFDである。そして「PIDの0でない素イデアルは極大イデアル」であることによって, どちらもPIDではないことが分かる($(x) \subset \mb{Z}[x]$, $(x)\subset \mb{R}[x,y]$はともに極大でない素イデアルである)。なお, 単項でないイデアルの存在を示してもよい。たとえば, $(x,2)\subset \mb{Z}[x]$や$(x,y) \subset \mb{R}[x,y]$などがある。\\
(2): $k[x_1,\cdots]$\\
(3): $\mb{Z}$. \\
(4): \bolm{\mb{Z}[\sqrt{-5}]}. 仮にこれがUFDだとすると\tf{既約元は素元になる.}よって既約元分解が一意である.  この環は整数環なので$\mb{Z}$の$\mb{Q}(\sqrt{-5})$における整閉包であり整閉整域である. ノルムを見ることで2,3は既約元であることが分かる. $1\pm \sqrt{-5}$も素元である. しかし
\[2\cdot 3 = (1+\sqrt{-5})(1-\sqrt{-5})\]
で分解が一意でなく, 矛盾.\qed

\end{egans}





\begin{egprob} (\cite{aoyukie})
$\mb{C}[x,y,z]/(z^2 - xy)$が正規環であることを示せ。
\end{egprob}
\begin{egans}

\end{egans}





\begin{egprob}
$\maf{p} = (2+\sqrt{-1})$を$A:=\mb{Z}[\sqrt{-1}]$のイデアルとする。$A_{\maf{p}}$は$\mb{Z}_{(5)}$上整ではないことを示せ。
\end{egprob}
\begin{egans}
$\dfrac{1}{2-\sqrt{-1}} = \dfrac{2-\sqrt{-1}}{5}\in A_{\maf{p}}$である。これを$x$とおく。もし$x$が$\mb{Z}_{(5)}$係数の整等式を満たせば, 適当な5で割れない整数$m$によって$mx$が$\mb{Z}$係数の整等式を満たすようにできる。$\mb{Z}[\sqrt{-1}]$は$\mb{Q}(\sqrt{-1})$における$\mb{Z}$の整閉包だが, $mx\notin \mb{Z}[\sqrt{-1}]$なので矛盾である。\qed
\end{egans}
\begin{rem}
$A$は$\mb{Z}$上整である。$S=\mb{Z}-5\mb{Z}$とするとき,$S^{-1}A$は$S^{-1}\mb{Z}$上整である(局所化は整従属を保存するため\footnote{\cite{atiyah}の5章参照})。これは上の問題とは異なる。つまり, $A\subset B$が整拡大でも$A_{\maf{p}^{c}} \subset B_{\maf{p}}$が整拡大とは限らない。なお, $S^{-1}A$は言葉でいえば「本質的には分母に5が来ない$\mb{Q}(\sqrt{-1})$の元の「集合」だから解答する際「分母に5が来てしまうようなものを見つければ勝てそう」と思える。
\end{rem}


\begin{egprob} (2019代数I期末1)
整域$A$は環$R$の部分環とする。$S=A\setminus{\{ 0 \}}$とする。$R$が$A$加群として平坦であり,  $R$の$S$による局所化$S^{-1}R$が整域であるとするとき, $R$も整域であることを示せ。
\end{egprob}
\begin{egans}
自然な準同型$\pi: R\to S^{-1}R$を考える。$R$が整域でないとする。ある$R$の元$x,y\neq 0$が存在して$xy=0$である。このとき$\pi(xy) = xy/1 = 0$であるから, $S^{-1}R$が整域であることより$x/1 = 0$または$y/1 = 0$である。$x/1=0$として考える。ある$s\in S$が存在して$sx = 0$である。したがって$s$を掛けるという写像$\phi:R\to R$は$\ker{\phi}$に$x$を含むので単射ではない。

一方で, $s$を掛ける写像$\psi:A\to A$は, $s\neq 0$と$A$が整域であることから単射である。$R$は$A$加群として平坦であるから, テンソル関手$\tens_{A} R$は単射を保つ。したがって
\[\psi \tens \mr{id}_{R}:  A\tens_{A} R \to A\tens_{A} R\]
は単射である。ところが, $A\tens_{A} R \cong R$により$\psi$は$\phi$そのものと同一視できるため, $\phi$が単射でないことに反する。よって矛盾であり$R$は整域である。
\end{egans}


\newpage
\begin{egprob} (数学II)
$p$を素数とし, 位数$p^2$の(可換とは限らない)環$A$を考える。
\ben
\item $A$が可換環であることを示せ。
\item $A$の同型類をすべて答えよ。
\een
\end{egprob}






\clearpage
%%%%%%%%%%%%%%%%%%%%%%%%%%%%%%%%%%%%%%%%%%%%%%
\subsection{演習問題}
\subsubsection{Level A}


\begin{prob}
2変数多項式環$\mb{Z}[t,s]$の部分環$\mb{Z}[t^2,ts,s^2]$は$\mb{Z}[x,y,z]/(y^2-xz)$に同型であることを示せ。
\end{prob}

\begin{prob}
$\mb{Z}[\sqrt{-5}]/(3) \cong \mb{F}_{3}\times \mb{F}_{3}$であることを示せ。
\end{prob}

\begin{prob} (H21留学生入試) [解答あり] \label{prob:H21F3}
Let $\mb{C}[x]$ be the polynomial ring over the field $\mb{C}$ of complex numbers.  We define a subset $I$ of $\mb{C}[x]$ by
\[I= \{ f(x)\in \mb{C}[x]\mid f(1) = f'(1) = 0 \},\]
where $f'$ denotes the derivative of $f$.  Show that $I$ is an ideal of $\mb{C}[x]$ and compute the dimension of the quotient ring $\mb{C}[x]/I$ as a vector space over $\mb{C}$.
\end{prob}

\subsubsection{Level B}


\begin{prob} (H19-I-3) \label{prob:H19I3}
$\mb{Z}[\sqrt{-3}]$のイデアルで$4$を含むものをすべて求めよ。
\end{prob}


\begin{prob} (\cite{aoyukie}, Ex2.7.7)
$\mb{C}[x,y,z]/(x^2+y^3+z^5)$が正規環であることを示せ。
\end{prob}

\begin{prob}
$A=\mb{Z}[\sqrt{-5}]$のイデアル$I=(5+4\sqrt{-5}, 15-3\sqrt{-5})$の素イデアル分解を求めよ。
\end{prob}

\begin{prob}
$A=\mb{Q}[x,y,z]/(x^2y^2-z^3)$とし, $p\in \mb{Q}[x,y,z]$の同値類を$\ovl{p}$とする。
\ben
\item $A$は整域であることを示せ。
\item $I=(\ovl{x},\ovl{y})$, $J=\ovl{x},\ovl{z})$はそれぞれ素イデアルであるか, 説明せよ。
\item 単項イデアル$(\ovl{x}\ovl{z})$を含む素イデアルで, 極大イデアルでないものを全て求めよ。
\een
\end{prob}

\begin{prob} (\tf{}) \checkbox\\
次の環の素イデアルの個数を求めよ。
{\small
\[\parena{\mb{Z}/2\mb{Z}\times \mb{Z}/4\mb{Z}\times \mb{Z}/12\mb{Z}}\tens_{\mb{Z}} \parena{\mb{Z}/14\mb{Z}\times \mb{Z}/42\mb{Z}\times \mb{Z}/84\mb{Z}\times\mb{Z}/420\mb{Z}} \]
}
\end{prob}

\begin{prob}
可換環の二つの単項イデアルの共通部分はまた単項イデアルですか?
\end{prob}

\subsubsection{Level C}

\subsubsection{Level X}
\begin{prob} (PIDだがEuclid整域でない例, $\mac{O}_{\mb{Q}(\sqrt{-19})}$)はPIDだがEuclid整域でないことを示せ。
\end{prob}




\subsubsection{Hint・Answer}
{\small
\begin{probans}
$\phi:\mb{Z}[x,y,z] \to \mb{Z}[t,s]$を$x\mapsto t^2,y\mapsto ts, z\mapsto s^2$で核を計算せよ.
\end{probans}
\begin{probans}
$\mb{Z}[\sqrt{-5}]\cong \mb{Z}[x]/(x^2+5)$で, $\mb{F}_{3}[x]/(x^2+5)$を考える.
\end{probans}
\begin{probans}
$I=((x-1)^2)$を示す。ただし右辺は$(x-1)^2$で単項生成されるイデアル。$I\supset ((x-1)^2)$は自明。$f(x)\in I$なら, $f(1)=0$から$f(x) = (x-1)g(x)$と書ける。$f'(x) = g(x) + (x-1)g'(x)$なので$f'(1) = g(1) = 0$. よって$g(x)=(x-1)h(x)$と書ける。よって$f(x) = (x-1)^2 h(x) \in ((x-1)^2)$.  $x=1$中心のTaylor展開により$\mb{C}[x]/I$は$[a(x-1)+b]$ ($a,b\in \mb{C}$)全体の集合である([]で剰余類を表す)。$[a(x-1)+b] = a[x-1] + b[1]$である。$a[x-1]+b[1] = c[x-1] + d[1]$ならば$[(a-c)(x-1) + (b-d)] = [0]$であるから, $(a-c)(x-1) + (b-d) \in I$である。$F(x) = (a-c)(x-1)+(b-d)$とおく。$F(1) = F'(1) = 0$からただちに$a=c, b=d$が従う。よって$\{ [1], [x-1] \}$は$\mb{C}[x]/I$を$\mb{C}$上生成し, かつ一次独立である。よって$\dim_{\mb{C}}\mb{C}[x]/I = 2$.
\end{probans}

%%%B

\begin{probans}
$\mb{Z}[\sqrt{-3}] \cong (\mb{Z}/4\mb{Z})[x]/(x^2-1)$を示し, そのイデアルを求めよ。この環の位数は16で, イデアルの位数は$16$である。この環の元の単元を決定せよ(8個ある)。
\end{probans}
\begin{probans}

\end{probans}
\begin{probans}

\end{probans}
\begin{probans}
素イデアル分解に現れる素イデアルは必ず$I$を含むから, $A/I$の素イデアルをすべて決定するとよい。$A\cong \mb{F}_{5}[x]/(x^2+5)$である。$A/I$は何と同型か?
\end{probans}
\begin{probans}
(1): ごり押しはなるべくしたくない。\textbf{\bolm{\phi(f(x,y,z)) \to f(t^3,s^3,t^2s^2)}という準同型$\mb{Q}[x,y,z]\to \mb{Q}[s,t]$を考えよ。$\ker{\phi}$は何か?}ところで, この写像の思い付き方は次のような思考過程による。$x^2y^2-z^3=0$で表される「図形」をある意味でパラメータづけたい。一変数では難しそうだから2変数$s,t$を使って考えてみる。与えられた$x,y$に対して$x=s^3, y=t^3$としてみれば$z$は$s^2t^2$であるべきだろう。(2): $A/I,A/J$は何と同型か。(3):$A/(\ovl{x}\ovl{z})$は何と同型か。そして, $\mb{Q}[x,y,z]$のどのような素イデアルと対応するのか。剰余環のイデアル対応定理は, 極大性を保存することにも注意せよ。
\end{probans}
\begin{probans}
$k[t^2,t^3]$で探せる.
\end{probans}
}


\clearpage



\section{イデアル}
環の\tf{イデアル}とは, 環$A$の加法に関する部分群であって, $x\in A$, $r\in I$ならば$rx\in I$を満たすようなものであった。とくに$A$でないイデアルのなかで極大なものを\tf{極大イデアル}といい, $x,y\notin \maf{p}\implies xy\notin\maf{p}$を満たすイデアル$\maf{p}$を素イデアルというのであった。また, 極大イデアルは素イデアルである。環準同型$\phi:A\to B$に対して$\ker{\phi}$は$A$のイデアルであり, $\im{\phi}$は$B$の部分環である。院試においては, 素イデアルを求める問題であったり, 環の同型を使って素イデアルであることを判定する問題が頻出である。




\subsection{準同型定理}

\begin{egprob}
$(x^2-y^3)$は$\mb{C}[x,y]$の素イデアルであることを示せ。
\end{egprob}
\begin{screen}
これは定義通りにやるようなものではない。$\mb{C}[x,y]/(x^2-y^3)$がある整域と同型であればよい。そのためには整域$A$として環準同型$\phi:\mb{C}[x,y]\to A$で$\ker{\phi}=(x^2-y^3)$となるものが欲しいが...
\end{screen}
\begin{egans}
曲線$y^3=x^2$の媒介変数表示を与えるということを考え, $\phi:\mb{C}[x,y] \to \mb{C}[t]$; $\phi(f(x,y)) = f(t^3,t^2)$とおく. このとき$\phi$は整域への環準同型であって, $\im{\phi}$は整域の部分環だから整域である. よって$\ker{\phi}$は素イデアルである. そして, $\ker{\phi} = (x^2-y^3)$を示す. [証明] $\ker{\phi}\subset (x^2-y^3)$を示す. $f(x,y)\in \ker{\phi}$であるなら,
\[f(x,y) = (x^2-y^3)Q(x)  +  xf_1(y) + f_2(y)\]
と書けるので,$\phi$を通し
\[t^3f_1(t^2)+f_2(t^2) = 0\]
である. 左辺の第1項はすべて奇数次で, 第2項はすべて偶数次なので, $\mb{C}[t] = \oplus_{d\geq 0} \mb{C}t^d$により$f_1(t^2),f_2(t^2) = 0$でなければならない. ここから$\mb{C}[x,y]$の元として$f_1,f_2 = 0$が従うからこの包含関係が示せた. 逆の包含は当たり前である. \qed
\end{egans}




\subsection{($\heartsuit$)剰余環の計算}
剰余環はなるべく簡単なものにして直すべきである。
\begin{prop} (剰余環まとめ)
\ben
\item (\tf{中国剰余定理}) $I,J$が互いに素なイデアル($I+J=A$)のとき, $IJ=I\cap J$であり, 次の環同型が存在する。
\[A/IJ \cong A/I \times A/J\]
\item $\{ A/I \text{のイデアル}\}$と $\{ A\text{のイデアルで}I\text{を含むもの} \}$には全単射がある。
\een
\end{prop}
\prf (1): 方針は$\pi:A\to A/I\times A/J$という自然な射から準同型定理を用いることである.\\
\fbox{$IJ=I\cap J$} $\subset$は常に成り立つ. $x\in I\cap J$とせよ. $1=i+j$となる$i\in I, j\in J$が存在し, $x = ix + jx \in IJ$なので示せた. \\
\fbox{$\pi$の核と像} これの核は明らかに$I\cap J$だが, 互いに素であることから$IJ=I\cap J$である. また全射でもある. 実際, $(x\bmod{I}, y\bmod{J})$を与えるとき, \bolm{\alpha=jx + iy}を考えると$\alpha \equiv jx \equiv x\pmod{I}$, $\alpha\equiv iy \equiv y\pmod{J}$となる. よって準同型定理から同型が従う. \\
(2):

\qed \\




\begin{prac} (中国剰余定理は無限個には拡張できない)\\
$\mathbb{Z}$と$\prod_{p:\mr{prime}} \mathbb{F}_{p}$は同型だろうか?
\end{prac}

\begin{egprob} (剰余環計算例1)
次の同型を示せ。
\ben
\item $\mathbb{C}[x]/(x^3+3x^2+2x) \cong \mathbb{C}^3$
\item $a^2-4b>0, =0, <0$のときそれぞれ$\mathbb{R}[x]/(x^2+ax+b) \cong \mathbb{R}^2, \mathbb{R}[\epsilon]/(\epsilon^2), \mb{C}$
\item $\mathbb{C}[x]/(x^4+1)$, $\mathbb{R}[x]/(x^4+1)$, $\mathbb{Q}[x]/(x^4+1)$
\een
\end{egprob}
\begin{point}
\ben
\item より一般に$\mathbb{C}[x]/(f(x))$の形を見たら, $f(x)$を因数分解しよう。$\mb{C}$なら一次式に分解できる。各一次式は極大イデアルを生成するが, 異なる極大イデアルは互いに素である。だから, もし$f(x)$が重根を持たないならこの剰余環は$\mathbb{C}^{\deg{f}}$と同型のはずだ。
\item $\mb{R}$係数なら, $\mb{C}\supset \mb{R}$という体拡大を視野に入れておこう。
\een
\end{point}
\begin{egans}

\end{egans}
より整数論的な問題も考えてみよう。「$\mb{Z}$の素数$p$が$\mb{Z}[i]$ではもはや素数でなくなってしまことがある」という話はよく聞く話であろう。これは代数的整数論の文脈ではよく調べられる対象であり, Dedekind環の拡大の「分岐・不分岐」という概念に対応する。

\begin{egprob}
\[\mathbb{Z}[\sqrt{-1}]/(p) \cong
\begin{cases}
\mathbb{Z}/4\mathbb{Z}\quad (p=2)\\
\mathbb{F}_{p}\quad (p\equiv 1\pmod{4})\\
\mb{F}_{p^2} \quad (p\equiv 3\pmod{4})
\end{cases}\]
\end{egprob}









\subsection{素イデアル・極大イデアル}
素イデアルは環における素数の役割を果たしており, 重要である。

\begin{prop} \checkbox (\tf{素イデアルと極大イデアル1}) $A$を環とする。
\ben
\item $\maf{p}$が素イデアル$\iff$ $A/\maf{p}$が整域
\item $\maf{m}$が極大イデアル$\iff$ $A/\maf{m}$が体
\item 極大イデアルは素イデアルである。
\item PIDの0でない素イデアルは極大である。とくに, 体$k$上の一変数多項式環$k[T]$の素イデアルは$(0)$であるか, ある既約多項式によって生成される。
\item  $I$を$A$のイデアルとするとき, $A/I$のイデアルは$A$のイデアルで$I$を含むものに全単射対応がある($\maf{p}\leftrightarrow \maf{p}/I$)。
\item
\een
\end{prop}
以下も可換環論ではよく使われる。
\begin{prop} (\tf{素イデアルと極大イデアル2})
\ben
\item $I$を真のイデアルとするとき, $I\subset \maf{m}$となる極大イデアル$\maf{m}$が存在する。
\item ($\spadesuit$) $S$を$A$の積閉集合とするとき, $S^{-1}A$の素イデアルは, $A$の素イデアルで$S$と交わらないものとの全単射対応がある( $\maf{p}\leftrightarrow S^{-1}\maf{p}$)。とくに, $\maf{p}$が素イデアルであるとき, $A_{\maf{p}}$の素イデアルは$\maf{p}$に含まれる素イデアルとも全単射対応がある。
\item \tf{環準同型$f:A\to B$と$\maf{q}\in \spec{B}$に対して, $f^{-1}(\maf{q})\in \spec{A}$である。}以後この写像$\maf{q}\mapsto f^{-1}(\maf{q})$を$\spec{f}$と書く。
\item ($\spadesuit$) $f:A\to B$により$B$を$A$代数と考える。$B$が$A$上整であるなら,\tf{ $\spec{f}$は極大イデアルを極大イデアルに移し, かつ極大でない素イデアルを極大でない素イデアルに移す. }
\item
\een
\end{prop}
\prf (1): $A\neq 0$なら極大イデアルが存在する. よって, $I\neq A$なら$A/I\neq 0$なのでこの環に極大イデアルが存在する. あとは剰余環の対応定理.\\
(2): 示すべきことは次である: $S^{-1}\maf{p}$が素イデアルであること, $\pi^{-1}(S^{-1}\maf{p})$が$S$と交わらない$A$の素イデアルであること, $\maf{p}\mapsto S^{-1}\maf{p}\mapsto \maf{p}$と移ること, $P$が$S^{-1}A$の素イデアルのとき $P\mapsto \pi{-1}(P) \mapsto P$と移ること. 読者に任せる. \\
(3): 素イデアルの定義にしたがってやれば簡単. \\
(4): $A\subset B$が整拡大のとき, $A$が体であることと$B$が体であることが同値. $A/f^{-1}(\maf{p}) \to B/\maf{p}$も整拡大であり, この二つを組み合わせて導ける. \qed





\begin{eg}
\ben
\item $\mb{C}[x]/(f(x))$ ($f\neq 0$)という環の素イデアルは, $f(x)$を割り切る任意の一次式$x-\alpha \bmod{f(x)}$で生成されるイデアルである。代数幾何ではこれを, 複素直線$\mb{C}$の点$\alpha$と同一視してしまう\footnote{($\spadesuit$)より一般に$A/I$という環の素イデアルは,$\spec{A}$の閉部分集合$V(I)$の点と同一視される。}
\item $\maf{p}\in \spec{A}$のとき, ある自然数$n$に対して$x^n \in \maf{p}$ならば$x\in \maf{p}$である。
\item 一般に$f:A\to B$が環準同型, $\maf{n}\in \spec{B}$が極大であっても, $f^{-1}(\maf{n})$が極大とは限らない。たとえば$A=\mb{Z}, B=\mb{Q}, f:\text{包含}, \maf{n} = (0)$とせよ。
\een
\end{eg}


\begin{egprob} \checkbox (\cite{midoriyukie} 補題8.11.12,13)
$p$を素数とし, $\zeta := \zeta_p$とおく($1$の$p$乗根)。
$\maf{p} = (1-\zeta_{p})\subset \mb{Z}[\zeta_{p}]$は素イデアルであることを示せ。また, $\maf{p}^{p-1} = (p)$を示せ。
\end{egprob}
\begin{screen}
剰余環が整域かどうかを見よう。$\mb{R}[x]/(x^2+1) \cong \mb{C} = \mb{R}[\sqrt{-1}]$のように, 代数的な元を添加した環は多項式環の剰余環で書くと良い。環同型の計算はよくやるものである。
\end{screen}
\begin{egans}
$\mb{Z}[\zeta]/(1-\zeta) \cong \mb{Z}[x]/(x^{p-1} + \cdots + 1)$であるから,
\[\mb{Z}[\zeta]/(1-\zeta) \cong \mb{Z}[x]/(x-1,x^{p-1}+\cdots + x + 1) \cong \mb{Z}[x]/(x-1,p) \cong h\mb{F}_{p}\]
これは体だから$\maf{p}$は素イデアル(極大イデアル)である。
\[x^{p-1} + x^{p-2} + \cdots + x + 1 = \disp\prod_{i=1}^{p-1}(x-\zeta^{i})\]
において$x=1$を代入すると
\[p = \disp\prod_{i=1}^{p-1} (1-\zeta^{i})\]
となる。ここで, $1-\zeta^{i}$は本質的に$1-\zeta$と同じである。すなわち, 同伴(単元倍程度の違いしかない)である。なぜなら
\[\dfrac{1-\zeta^{i}}{1-\zeta} = \zeta^{i-1}+\cdots + \zeta + 1\]
かつ, $ij\equiv 1 \pmod{p}$となる$j>0$を取ったとき
\[\dfrac{1-\zeta}{1-\zeta^{i}} = \dfrac{1-\zeta^{ij}}{1-\zeta^{i}} = 1 + \zeta^{i} + \cdots \in \mb{Z}[\zeta]\]
となるからである。よって$\maf{p} = (1-\zeta^{i})$であるからイデアルの等式
\[(p) = \maf{p}^{p-1}\]
を得る。\qed
\end{egans}



\begin{egprob} (Quiz, by Li)\\
$k$を代数閉体とする。$x,y$を不定元とし,2変数多項式環の剰余環 $A=k[x,y]/(xy)$を考える。$A$の素イデアルを決定せよ。
\end{egprob}

\begin{egans}
$\maf{p}\in \spec{A}$とする。$\maf{p}$は$k[x,y]$の素イデアルで$(xy)$を含んでいるものと思える。$xy\in \maf{p}$だから$x\in \maf{p}$か$y\in \maf{p}$が成り立つ。議論は同様なので$x\in \maf{p}$として考える。このとき$\maf{p}$は$k[x,y]/(x)\cong k[y]$の素イデアルに対応する。$k$は代数閉体だから任意の次数が2以上の多項式は可約である。また, 一次式$y-\alpha$ ($\alpha\in k$)はすべて既約である。よって$k[y]$の素イデアルは$(0)$か$(y-\alpha)$であり,これらは $\maf{p} = (x)/(xy), (x,y-\alpha)/(xy)$に対応する。$y\in \maf{p}$の場合も同様に$\maf{p}=(y)/(xy), (x-\alpha, y)/(xy)$に対応する。以上より,
\[(x)/(xy), (y)/(xy), (x,y-\alpha))/(xy), (x-\alpha, y)/(xy)\]
である。
\end{egans}
\begin{rem}
この中で極大イデアルは後ろ二つである。この極大イデアルは座標平面($k\times k$)の$xy=0$で定義される図形上の$(0,\alpha)$という点に対応していると思える。(前二つの素イデアルのほうは, それぞれ$y$軸 ($x=0$), $x$軸 ($y=0$) を表していると思える。)このようにして, $xy=0$上の点と$\spec{A}$の極大イデアルが1:1に対応する(\tf{Hilbeltの零点定理})。

素イデアルの集合を$\spec{A}$と書いたが,とくに$A$が有限生成$k$代数(つまり適当なイデアル$I$と自然数$n$によって$k[x_1,\cdots, x_n]/I$と同型なもの)であるときの$\spec{A}$は代数幾何・スキーム論の大きな動機になっている。つい前に述べた話はより一般的な現象であって,すなわち$k$が代数閉体で, $A= k[x_1,\cdots,x_n]/I$は有限生成$k$代数であるときに, $\spec{A}$の閉点(=極大イデアル)は, $I$に属す多項式の共通零点$Z(I) = \{ (s,t)\in k^2 \mid f(s,t)= 0 \quad (\forall f\in I) \}$と1:1対応がある, ということである。
\end{rem}


\begin{eg}
$\spec{\mb{Z}[x]}$を決定しよう. $\maf{p}\in \spec{\mb{Z}[x]}$とする. \\
\step{1}{ある素数$p$が$\maf{p}$に属す場合} このような素数$p$は一意である. $\maf{p}$は, $\mb{Z}[x]/p\mb{Z}[x]$の素イデアルに対応している.
\[\mb{Z}[x]/p\mb{Z}[x] \cong \mb{F}_{p}[x]\]
なので, $\maf{p}$は右辺の素イデアルに対応する. 右辺はPIDだから, その素イデアルは$(0)$か$(f)$ ($f$は既約多項式)である. $\witi{f} \bmod{p} = f$であるような$\witi{f} \in \mb{Z}[x]$を任意に取ると, これも既約である(もし可約ならそれを$\bmod{p}$して$f$も可約になってしまうから). これらをこの同型で引き戻すことで
\[p\mb{Z}[x]/p\mb{Z}[x],\quad (p,\witi{f})\mb{Z}[x]/p\mb{Z}[x]\]
に対応し, $\spec{\mb{Z}[x]}$の中の
\[p\mb{Z}[x], \quad (p,\witi{f}(x))\mb{Z}[x]\]
という素イデアルに対応している. \\
\step{2}{$\maf{p}\cap\mb{Z} = \set{0}$のとき} $S=\mb{Z}-\set{0}$は積閉集合である.  $S^{-1}\mb{Z}[x]$の素イデアルと$\mb{Z}[x]$の素イデアルで$S$と交わらないものが自然に対応し, $\maf{p}\cap \mb{Z} = \varnothing$なので, $S^{-1}\maf{p} \in \spec{S^{-1}\mb{Z}[x]}$である. ここで
\[S^{-1}\mb{Z}[x] \cong \mb{Q}[x]\]
であり, $\mb{Q}[x]$はPIDだから$S^{-1}\maf{p} = (0), (f)$ ($f$は$\mb{Q}$上既約な多項式) である. これを局所化写像$\mb{Z}[x]\to S^{-1}\mb{Z}[x]$で引き戻して
\[(0), (f)\]
という素イデアルを得る. 逆にこれらの形のものが必ず素イデアルになるということも示せる.
\end{eg}
\begin{prac}
$S^{-1}\mb{Z}[x]\cong \mb{Q}[x]$を具体的に証明せよ. (\tf{Hint}: 自然な包含$\mb{Z}[x]\to \mb{Q}[x]$は$S$を可逆にするので, 普遍性によって$S^{-1}\mb{Z}[x]\to \mb{Q}[x]$を構成できる. これの逆対応を自然に与えよ. )
\end{prac}



























次に零点定理を$\mb{Z}$上で考えることに関連する問題をひとつ考える. これも大変難しいが, 知っていると便利である.

\begin{egprob} ($\spadesuit\spadesuit$)
有限生成$\mb{Z}$代数$A=\mb{Z}[t_1,\cdots, t_n]$を考える. $A$の極大イデアルは必ず素数$p$を含むことを証明せよ.
\end{egprob}
\begin{egans}
$\maf{m}$が素数$p$を含まないとせよ. $S=\mb{Z}-\set{0}$は積閉集合である. $B=S^{-1}A$とすると$B=\mb{Q}[t_1,\cdots,t_n]$である. $\maf{m}\in \mr{mSpec}(A)$は$(S^{-1}\maf{m} =) \maf{m}B \in \mathrm{mSpec}(B)$に対応する. 体$k=A/\maf{m}$は仮定より有理数体$\mb{Q}$を含むため, $S$の元はすでに可逆である. よって$S^{-1}k = k = S^{-1}A/S^{-1}\maf{m} = B/\maf{m}B$であり, $B/\maf{m}B$は有限生成$\mb{Q}$代数$B$の剰余体だから\tf{Zariskiの補題}より$B/\maf{m}B$は$\mb{Q}$の有限次拡大である. つまり$A/\maf{m} \cong B/\maf{m}B$の任意の元は$\mb{Q}$上代数的. $\ovl{t_i} = t_i + \maf{m}$と書くと, $\ovl{t_i}$は$\mb{Q}$上代数的なので, その$\mb{Q}$上最小多項式を$m_i(X)\in \mb{Q}[X]$と表そう. $\mb{Z}$上で, $m_i(X)$のすべての係数で生成される$\mb{Z}$代数$C$を考えると($i$が\tf{有限個だから！})これは有限個の有理数$p_j$ ($1\leq j\leq N$)によって
\[C=\mb{Z}[p_1,\cdots, p_N]\]
と書けるわけだが, $m_i(X)\in C[X]$である. よって$\ovl{t_i}$は$C$上整である. $i$は任意だから,
\[C = \mb{Z}[p_1,\cdots, p_N] \subset \mb{Z}[\ovl{t_1},\cdots, \ovl{t_n}] = A/\maf{m}\]
は\tf{整拡大.} $A/\maf{m}$は体なので, \tf{$C$も体.} すなわち$C=\mb{Q}$になるしかないが, $p_i$を既約分数にしておいて, その分母に現れない素数$q$を取って$\frac{1}{q}$を考えると$\frac{1}{q} \notin \mb{Z}[p_1,\cdots, p_n]$が容易にわかるので矛盾.\footnote{この部分はまあ非常に簡単だが, $\mb{Q}$が有限生成$\mb{Z}$代数にならないという問題と同じ要領である.}\qed
\end{egans}

次のような命題から上の問題の別の答案を作ることもできる.

\begin{prop} (\cite{atiyah}命題7.8)
環の列$A\subset B\subset C$を考え, $A$はNoether環,  $C$は有限生成$A$代数とする。\tf{$C$が有限生成$B$加群であるか, $C$が$B$上整であることが分かっているとする。}このとき, $B$は$A$加群として有限生成である。
\end{prop}

\begin{eg} (\cite{atiyah} Ex7.6)
$C=\mb{Z}[t_1,t_2,\cdots, t_n]$という有限生成$\mb{Z}$代数を考える. このとき, $C$は決して標数0の体にはならない. なぜなら, もしなれば標数が$0$なので$\mb{Z}\subset \mb{Q}\subset C$という部分環の包含があり, $C$は有限生成$\mb{Q}$代数であるような体だから, \tf{Zariskiの補題}よりそれは$\mb{Q}$上の有限次拡大体であって, とくに有限生成$\mb{Q}$加群である.  これを先の命題に適用すると$\mb{Q}$が$\mb{Z}$加群として有限生成であるという結論が得られ, 矛盾である. よって体であるならば有限体であるしかない. ここからさらに次のことが言える: \tf{$\mb{Z}[T_1,\cdots, T_n]$の極大イデアル$\maf{m}$は必ずある素数$p$を含む.}[証明] $\mb{Z}[T_1,\cdots, T_n]/\maf{m}$は$\mb{Z}$代数として有限生成な体なので, 先に示したことから標数は$p>0$である. よって$p\in \maf{m}$が従う. \qed
\end{eg}


\newpage

\subsection{演習問題}

\subsubsection{Level A}
\begin{prob}
$\maf{p}\in \spec{A}$のとき, $\maf{p}[x]\in \spec{A[x]}$であることを示せ。ただし$\maf{p}[x]$とは, 係数が$\maf{p}$の元であるような$A[x]$の元全体である. また, $\maf{p}[x]\cap A = \maf{p}$を示せ.
\end{prob}


\subsubsection{Level B}
\begin{prob}
$k$を体とする。$k[T_1\cdots, T_n]$の極大イデアルは$n$元で生成されることを示せ。
\end{prob}



\subsubsection{Level C}
\begin{prob}
$n$を2以上の整数とし, $\zeta = \exp{2\pi \set{-1} n}$を$1$の原始$n$乗根とする.
\[R = \set{f(X,Y) \in \mb{C}[X,Y] \mid f(\zeta X, \zeta Y) = f(X,Y)}\]
とおく. 以下の問に答えよ.
\ben
\item $\mb{C}$代数として, $R$は$n+1$個の元$X^n,X^{n-1}Y,\dots, Y^n$で生成されることを示せ、
\item 複素数$a,b,c,d$に対し$m_{a,b}=(X-a,Y-b)$, $m_{c,d} = (X-c,Y-d)$を$\mb{C}[X,Y]$のイデアルとする. $m_{a,b}\cap R = m_{c,d} \cap R$が成り立つための, $a,b,c,d$に関する必要十分条件を求めよ.
\een

\end{prob}

\begin{prob} ($\bigstar 8$) \label{prob:atimac_prop7.8_apply}
\ben
\item $\mb{Z}$上有限生成加群な体$K$は有限体であることを示せ。(\cite{atiyah} 演習7.4)
\item 有限生成$\mb{Z}$代数$A$の極大イデアル$\maf{m}$に対し, $A/\maf{m}$は有限体であることを示せ。とくに, $\maf{m}\cap \mb{Z} = p\mb{Z}$となる素数$p$が存在することを示せ。(\cite{gotou}, 演習9.3)
\een
\end{prob}

\begin{prob} ($\bigstar 10$, H17数学II問1)
$\mb{Z}[x_1,\cdots, x_n]$の極大イデアルは$n+1$個の元で生成されることを示せ。
\end{prob}


\subsubsection{Hint・Answer}
{\small
%=================Lev A

\begin{probans}
$A[x]/\maf{p}[x] \cong (A/\maf{p})[x]$を準同型定理から示せ. 後半は集合論的に素朴にやる.
\end{probans}


%================================Level B
\begin{probans}
帰納法で示す. $n=1$のときはPIDだから明らか. $n\geq 2$とし,  $n-1$での成立を仮定. $M$を$k[T_1,\dots, T_n]$の極大イデアルとする. 単射準同型$i:k[T_1,\dots, T_{n-1}]\to k[T_1,\dots, T_n]$を考える. 例題\ref{egprob:fg-k-alg-closedpt-to-closedpt}により$i^{-1}(M)$は極大である. よって$i^{-1}(M)$は$n-1$個の元で生成される. $k_1 := k[T_1\dots, T_{n-1}]/i^{-1}(M)$とおくとこれは体である. $i^{-1}(M)[T_n] = i^{-1}(M)k[T_1,\dots, T_n]$なので, $k[T_1,\dots T_n]/i^{-1}(M)[T_n] \cong k_1[T_n]$はPIDである. $M$の左辺の環における像もまた極大イデアルなので, それは$k_1[T_n]$の単項イデアル$(F)$に対応する. $F$の係数は$k_1$の元であるが, その各係数から$k[T_1,\dots, T_{n-1}]$のリフトを持ってきて$F_0\in k[T_1,\dots, T_{n-1}][T_n]$を取ると, $F_0\in M$でありかつこれが$M/(i^{-1}(M)[T_n])$を生成する. よって$M = F_0k[T_1,\dots, T_n] + i^{-1}(M)k[T_1,\dots, T_n]$を得る. $i^{-1}(M)$の生成元は$n-1$個だから,帰納法が回る.
\end{probans}




%========================Level C
\begin{probans}
\end{probans}
\begin{probans}
\end{probans}
\begin{probans}
\end{probans}

}















\clearpage
\section{($\spadesuit$)局所化と整従属}
\subsection{局所化}
面接で聞かれてもよいように, 局所化の構成についておさえておこう. $S\subset A$を積閉集合で0は含まないとし, $S^{-1}A$を次の集合とする: $S\times A$の元$(s,a)$, $(s',a')$に対し, 次の関係を定める. すなわち, $(s,a)\sim (s',a') \iff \exists t\in S, t(sa' - s'a) = 0$. これは実は同値関係をなし, $(s,a)$の同値類を$\dfrac{a}{s}$と書く. 定義から, $sa' = as'$なら$\dfrac{a}{s} = \dfrac{a'}{s'}$である. \par
この構成も大事であるが, 次の\tf{普遍性}も理論的には需要であり, 写像を作るときに便利である.
\tbox{
$\phi:A\to B$を環準同型で, $\phi(S) \subset B^{\times}$なるものとする. このとき, $\phi$は次のような準同型$\ovl{\phi}:S^{-1}A\to B$を引き起こす.
\[\xymatrix{
A\ar[r]^{\phi}\ar[d] & B \\
S^{-1}A \ar[ur]_{\ovl{\phi}} &
}
\]
ただし, $A\to S^{-1}A$は$a\mapsto \dfrac{a}{1}$を満たす自然な準同型.
}{}



\begin{prop}
\ben
\item $A$を環で, $f$が零因子でないとき, $A_{f} \cong A[x]/(xf - 1)$.
\item 局所化は完全関手である. したがって, 局所化と, 剰余・直和・テンソルなどは交換できる.
\item (\tf{重要})$S$を積閉集合とする. $X=\spec{A}$の部分集合$S^{-1}X:=\set{\maf{p}\in \spec{A}\mid \maf{p}\cap S = \varnothing}$と$\spec{S^{-1}A}$の間に自然な全単射がある.
\een
\end{prop}
\prf
(1): $F:A[x] \to A_f$を$x\mapsto 1/f$なる$A$代数の射とする. これは$F(xf - 1) = 0$を満たすので$\witi{F}:A[x]/(xf-1) \to A_{f}$が誘導される. 逆に$G:A\to A[x]/(xf-1)$を自然な$A$代数の射とするとき, $G(f^n) = f^n + (xf-1)$は$x^n + (xf-1)$の逆元である. よって普遍性から$\witi{G}:A_{f} \to A[x]/(xf-1)$が引き起こされる. $\witi{G}\circ \witi{F}$と$\witi{F}\circ \witi{G}$が恒等写像であることは容易に分かる.
\qed

\begin{rem}
$A_{f} = 0 \iff f$が冪零元である. このようなとき, $A[x]/(xf-1)$も0である. 実際, $xf - 1$が$A[x]$の単元であることを示せばよい. 実際, $f^n = 0$なら, $\sum_{j=0}^{n-1} (xf)^{j}$が$xf-1$の逆元. \cite{atimac}のEx1.2も参照せよ.
\end{rem}


\begin{egprob}
環$A$に対し$D(f) = \set{\maf{p}\in \spec{A} \mid f\notin \maf{p}}$とおく.
\ben
\item $D(f) \subset D(g)$ならば$f^n = gr$となる$n\in \mb{N}$, $r\in A$が存在することを示せ.
\item $D(f) \subset D(g)$とし, (1)のように$n,r$を取る. 環準同型$A_{g} \to A_{f}$で$\dfrac{1}{g} \mapsto \dfrac{r}{f^n}$となるものが存在することを示せ.  これは$r,n$に依らないことを示せ.
\item $D(f) = D(g)$ならば, $A_f\cong A_{g}$であることを示せ.
\een
\end{egprob}
\begin{egans}
(1): $V(fA) \supset V(gA)$により$g$を含む素イデアルは$f$も含むので, $g\in \bigcap_{\maf{p}\ni g} \maf{p} = \sqrt{gA}$が従う. ここから直ちに得られる.\\
(2): $f^{n_1} = gr_1$, $f^{n_2} = gr_{2}$とせよ. 示したいことは$A_{f}$上で$r_1/f^{n_1} = r_2/f^{n_2}$となることである. これは, ある$f^k$が存在して$f^k(r_1f^{n_2} - r_2f^{n_1}) = 0$となることと同値だが, これは$r_2f^{n_1} = gr_1r_2 = r_1f^{n_2}$により明らかである. \\
(3):
\end{egans}
次は初見では案外難しい.
\begin{egprob}
次を示せ.
\[\sqrt{(0)} = \bigcap_{\maf{p}\in \spec{A}} \maf{p}\]
\end{egprob}
\begin{point}
$\supset$が難しい. 実のところ, これは選択公理を認めてやらねばならない. ポイントは, $\spec{A} \neq \varnothing \iff A\neq 0$を利用することである(この命題はZornの補題を用いるのであった).
\end{point}
\begin{egans}
$\subset$はやさしい(Check!). $\supset$を見る. $f$は右辺に含まれる元とし, 局所化$A_{f}$を考える. もし$f$が冪零元でなければ, $A_f\neq 0$である.
よって$\spec{A_f}$は空でないので, ここから素イデアル$P$が取れる. その$P$は$A$の$f$を含まない素イデアルに対応し, 矛盾である. もし局所化の対応が苦手であるなら, 次のようにせよ: $f$が冪零でなく右辺に含まれる元とする. \aka{$f$の冪を含まない}イデアル全体の集合$\Sigma$は$(0)$を持つので空でなく, Zornの補題の適用条件を満たすことが確認できる. $\Sigma$の極大元が少なくとも一つ存在する. そのイデアルを$I$とせよ. この$I$が実は素イデアルである. [証明] $x,y\notin I$で$xy\in I$として矛盾を導く. 極大性より$f^n\in I+(x), f^m\in I+(y)$なので$f^{n+m}\in I$だがこれは$I\in \Sigma$に反する.\qed
\end{egans}
これと似ているが重要な定理として, 次のようなものも習うだろう.
\begin{egprob}
$A$を有限生成$k$代数とするとき, 次を示せ.
\[\bigcap_{\maf{m}\in \mr{mSpec}A} \maf{m} = \sqrt{(0)}\]
\end{egprob}
\begin{egans}
$\supset$は明らか. 左辺から元$f$を取り, これが冪零ではないとする. このとき$A_{f}\neq 0$で, $A_{f}\cong A[T]/(Tf-1)$は有限生成$k$代数. $A_{f}$の極大イデアル$\maf{m}$を取ると, $A_{f}/\maf{m}$は$k$の有限次拡大(Zariskiの補題). $\pi:A\to A_{f}$を自然な写像とすると, 単射$A/\pi^{-1}(\maf{m}) \to A_{f}/\maf{m}$が引き起こされる. 構造射$k\to A_{f}/\maf{m}$は有限次拡大の埋め込みなので整準同型であり, したがって$A/\pi^{-1}(\maf{m})\to A_{f}/\maf{m}$もそうである. すると$A/\pi^{-1}(\maf{m})$は体でなければならないので, $\pi^{-1}(\maf{m})$は極大イデアルとなる. これは$f$を含まないが, それは$f$の取り方に反する. \qed
\end{egans}


\begin{egprob} (H 数学II,1)
単位元1を持つ可換環$A$のイデアル$I$について, $A$の元$a$がイデアル$I$の根基$\sqrt{I}$に含まれるための必要十分条件は, $A$上の一変数多項式環$A[T]$の中で$I$と$1-aT$が生成するイデアルが1を含むことであることを示せ.
\end{egprob}
\begin{point}
問題文は簡単に理解できよう. しかし, このままやるのではなく, 少し簡単にする方法を考えて混乱しないようにすることが環論ではよく重要になる. あるいは, 簡単だが本質的な場合に示すことでヒントを得るのである.
\end{point}

\begin{egans} $I=(0)$の場合を考える(実は剰余環を考えればこれに帰着される).
$a\in \sqrt{0}$のとき, $(1-aT) = A[T]$を示す. これは$1-aT$が単元であることに同値. $a^n = 0$なる$n$を取ると
\[(1-aT)\sum_{j=0}^{n-1}a^jT^j = 1- a^nT^n = 1\]
より単元. よって$(1-aT) = A[T]$. 逆に$(1-aT) = A[T]$であるとせよ. このとき$A[T]/(1-aT) \cong A_{a} = 0$なので$A_a$の中で$1/1 = 0/1$であり, これより$a^n = 0$という関係式が得られる. \par
一般の$I$の場合を考えよう. $IA[T]+(1-aT) = A[T]\iff (1-aT)\dfrac{A[T]}{IA[T]} = \dfrac{A[T]}{IA[T]} \iff 1-aT \bmod{IA[T]}$が$\dfrac{A[T]}{IA[T]} $で単元$\iff 1 - (a\bmod{I})T が(A/I)[T]$で単元$\iff a\bmod{I}\in A/Iが冪零 \iff a\in \sqrt{I}$.
\end{egans}


\subsection{整拡大}
\begin{prop}
環の列
$A\subset B\subset C$で$C$が$B$上整かつ$B$が$A$上整とするなら$C$は$A$上整.
\end{prop}
\prf
$x\in C$を取る. $x$の$B$係数整等式の係数$b_0,\dots, b_{n-1}$を取り, 環$A[b_0,\dots, b_{n-1}]$を考える. $x$は$A[b_0,\dots, b_{n-1}]$上整であるから, $A[b_0,\dots, b_{n-1},x]$は$A[b_0,\dots,b_{n-1}]$加群として有限生成. $b_{i}$は$A$上整なので, $A[b_0,\dots, b_{n-1}]$は$A$加群として有限生成. よって$R:=A[b_0,\dots, b_{n-1},x]$は$A$加群として有限生成. $R\supset A[x]$であるから, 主張は次の命題の(i)$\iff$(iii)から示される. \qed
\begin{prop} (\cite{atimac} 命題5.1. これは整拡大を理論的に簡単に取り扱うために重要な命題だと思う.) $A\subset B$を環の列とするとき, $x\in A$に対して次は同値.
\ben
\item $x$は$A$上整.
\item $A[x]$は$A$加群として有限生成.
\item $A$加群として有限生成である環$C\subset B$が存在して, $A[x]\subset C$.
\item $A$加群としては有限生成であるが, $A[x]$加群としては忠実である加群$M$が存在する.
\een
\end{prop}
\qed



\tbox{
\begin{theo}
$A\subset B$を整域の整拡大とするとき, $A$が体であることと$B$が体であることは同値である.
\end{theo}
}{}

\begin{prac} \label{prac:closedpoint_correspondence}
整拡大$A\subset B$において, $\maf{n}\subset B$が極大であることと$\maf{n}\cap A\subset A$が極大であることは同値であることを示せ.
\end{prac}

\begin{cor}
$A\subset B$を整拡大とする. $\maf{q_1},\maf{q}_2 \in\spec{B}$で$\maf{q}_1\subset \maf{q}_2$かつ$\maf{q}_1\cap A = \maf{q}_{2}\cap A$を満たすものとすると, $\maf{q}_1 = \maf{q}_{2}$となる.
\end{cor}
\begin{rem}
すなわち, 整拡大に対応する射$f:\spec{B}\to \spec{A}$の$\maf{p}\in \spec{A}$におけるファイバー$f^{-1}(\maf{p})\cong \spec{A_{\maf{p}}/\maf{p}A_{\maf{p}}}$の異なる2点には\tf{包含関係がない.} 気持ち的には ファイバーの2点$x,y\in f^{-1}(\maf{p})$はバラバラな印象がある.
\end{rem}
\prf この命題を見て, \tf{ファイバーのことを思い出すとよい}ので, 局所化して考えるのが妥当である. $\maf{p} = \maf{q}_{1}\cap A = \maf{q}_{2}\cap A$とする. $B$の整閉集合$A-\maf{p}$による局所化$B_{\maf{p}}$は$A_{\maf{p}}$上整である. $\maf{q}_{1}B_{\maf{p}}$, $\maf{q}_{2}B_{\maf{p}}$は$B_{\maf{p}}$の素イデアルであり, 仮定より$\maf{q}_{1}B_{\maf{p}} \subset \maf{q}_{2}B_{\maf{p}}$を満たす. $\maf{p}A_{\maf{p}} = \maf{q}_{1}B_{\maf{p}}\cap A_{\maf{p}} = \maf{q}_{2}B_{\maf{p}}\cap A_{\maf{p}}$となる. \ao{ $\maf{p}A_{\maf{p}}\subset A_{\maf{p}}$は極大なので, 練習問題により$\maf{q}_{i}B_{\maf{p}}$も極大である.} よって$\maf{q}_{i}B_{\maf{p}}$はともに極大で, 包含があるのでそれらは等しい. 局所化との素イデアル対応定理によって$\maf{q}_{1} = \maf{q}_{2}$でなければならない. \qed

{\small
\begin{rem} 全く同じ証明だが, 代数幾何的な書き方にするとこうなる.\\
$R_{P} \to A_{Q_2}$は単射整拡大. よって$\spec{A_{Q_2}}\to \spec{R_P}$で閉点同士が対応する(\ref{prac:closedpoint_correspondence}参照). $Q_2A_{Q_2}, Q_1A_{Q_2}$はこの写像で送ってともに$PR_P$という閉点になるので, $Q_iA_{Q_2}$は閉点. $A_{Q_2}$は局所環だから$Q_1A_{Q_2} = Q_2A_{Q_2}$である. これでは$Q_1=Q_2$となるので矛盾.\qed
\end{rem}
}

実は上の命題が過去問で登場しているが, 少し違う流れで示す問題になっている(演習問題参照). \par
試験で役立つ感じはしないが, 代数学Iで扱うことなので復習しておく.
\lefba{
\begin{theo} (\tf{上昇定理})
$A\subset B$を\tf{整拡大}, $\maf{p}_1\subset \dots \subset \maf{p}_n$を$A$の素イデアル列とする. \bolm{1\leq m<n}に対して, $\maf{q}_{i}\cap A = \maf{p}_{i}$ ($1\leq i\leq m$)となるような素イデアル列$\maf{q}_{1}\subset \dots \subset \maf{q}_{m}$がもし有れば, この昇鎖は
\[\maf{q}_{1}\subset \dots \subset \maf{q}_{m}\subset \maf{q}_{m+1} \subset \dots \subset \maf{q}_{n},\quad \maf{q}_{i}\cap A = \maf{p}_i,  (1\leq i\leq n)\]
に延長できる.
\end{theo}
}
\prf $n=2$, $m=1$の場合を示す. $A/\maf{p}_{1} \subset B/\maf{q}_{1}$であり, これは整拡大. よってLying over theoremより$B/\maf{q}_{1}$のある素イデアル$\maf{q}_2/\maf{q}_1$が$\maf{p}_{2}/\maf{p}_1$の上にあるのでこの$\maf{q}_{2}$を取れば良い. \qed
次は覚えておくとよい.
\begin{egprob} \label{egprob:dimA_dimB_integral_ext}
\ben
\item $\phi:A\to B$が整準同型であるとき, $\dim{B}\leq \dim{A}$であることを示せ.
\item $\phi:A\to B$が単射整準同型であるとき, $\dim{A} = \dim{B}$を示せ.
\een
\end{egprob}
\begin{egans}
(1): $\maf{q}_0\subsetneq \maf{q}_1$とする. $A$に引き戻して$\maf{p}_0 \subset \maf{p}_1$を得る. このとき$\psi:A/\maf{p}_0\to B/\maf{q}_0$は整準同型. $x\in \maf{q}_1$とする. $y\in \maf{q}_1\setminus{\maf{q}_0}$とおく. $y\bmod{\maf{q}_0}$は$A/\maf{p}_0$上整. その整等式で最小次数のものを考え, その定数項を$a_0$とすれば,  最小性から$a_0 \not 0 \in A/\maf{p}_0$であり, $a_0 = -x (\dots)$の形で書け, $x\neq 0$と最小性から$a_0\neq 0$, つまり$a_0$の代表元が$\maf{p}_0$に属していない.  $a_0\in \maf{q}_1/\maf{q}_0 \cap A/\maf{p}_0 = \maf{p}_1/\maf{p}_0$で$a_0\neq 0$だから$\maf{p}_0\subsetneq \maf{p}_1$が分かる. よって$\dim{B} \leq \dim{A}$が従う. \qed

(2): $\maf{p}_0\subsetneq \maf{p}_1$とする. $\maf{q}_0  =  \maf{p}_0^{c}$なる$\maf{q}_0$がある. 上昇定理のように, 剰余で$A/\maf{p}_0 \to B/\maf{q}_0$に落としてLying over Thmを考えれば$\maf{p}_1$の上にある$\maf{q}_1\supsetneq \maf{q}_0$が作れるから$\dim{A}\leq \dim{B}$が分かる.
\end{egans}





\lefba{
\begin{theo} (\tf{下降定理}) \\
$A\subset B$を\bolm{Aが整閉整域, Bが整域の整拡大}とするとき, 上昇定理で$\subset$を$\supset$に置き換えた主張が成り立つ.
\end{theo}
}
これのある種の反例が過去問に存在する.
\begin{egprob} (H27専門科目)
$A=\mb{Z}[X,Y]/(Y^2-6X^2)$, $B=\mb{Z}[X,T]/(T^2-6)$とおく. $A$における$X,Y$の像を$x,y$とし, $B$における$X,T$の像を$x',t$とする. また,
\[P_1=(x,y,5)A, P_2 = (x-y,5)A, Q_1 = (x',t+1)B)\]
とする.
\ben
\item 単射環準同型$\phi:A\to B$で$\phi(x) = x'$, $\phi(y) = x't$であるものが存在することを示せ.\footnote{これはブローアップ写像と関連すると思われますが, 詳細は未確認です. }
\item $P_1,P_2\in \spec{A}$で$P_2\subsetneq P_1$を示せ.
\item $\phi$により$A\subset B$とみなす. $Q_1\in \spec{B}$で, $Q_1\cap A = P_1$を示せ.
\item (\tf{下降定理の反例}) $Q_2\in \spec{B}$で$Q_2\subset Q_1$かつ$Q_2\cap A = P_2$であるものは存在しないことを示せ.
\een
\end{egprob}
\begin{egans}
(1): $X\mapsto X, Y\mapsto XT$なる環準同型を考えるとこれが $(Y^2-6X^2)$を$(T^2 - 6)$に移るので$\phi:A\to B$が引き起こされる. $A$の元は$a(x)y + b(x)$の形に書ける. $\phi$で移すと$a(x') + x'tb(x')$であり, これが0なら$a(x') = 0$, $b(x')=0$でなければならないので, $a(x) + yb(x) = 0$であり単射となる. \\
(2):
\bealn{
&A/P_1 \cong \mb{Z}[X,Y]/(Y^2-6X^2, X,Y,5) = \mb{Z}[X,Y]/(5,X,Y) \cong \mb{F}_{5}\\
&A/P_2 \cong \mb{Z}[X,Y]/(X-Y, 5) = \mb{Z}[X,Y]/(5,X-Y) = \mb{F}_{5}[X]
}
で両方整域なので$P_i\in \spec{A}$. $P_2\subset P_1$は明らかで, 剰余環が異なるので$P_1\neq P_2$. \\
(3): $Q_1\cap A = \phi^{-1}(Q_1)$である.  $\phi(x) = x' \in Q_1$, $\phi(y) = x't = x'(t+1) - x' \in Q_1$, $\phi(5) = 5 = t^2 - 1 = (t-1)(t+1)\in Q_1$なので$\phi^{-1}(Q_1) \supset Q_1$. 逆に$a(x)+yg(x) \in \phi^{-1}(Q_1)$なら$a(x') + x'tb(x') \in Q_1$だから$a(x') \in Q_1$なので$a(x)$の定数項の整数は$Q_1$に属す. $5 = t^2-1 \in Q_1$でこれが$Q_1$の持つただひとつの素数であるから, $a(x)$の定数項は5の倍数である. よって$a(x) + yg(x)\in (x,y,5) = P_1$が分かる. \par
(4): 仮にあるとしよう. このとき$5\in P_2 \subset Q_2$である. よって, $\mb{F}_{5}$代数の射
\[\mb{F}_{5}[X]\cong A/P_2 \hookrightarrow B/Q_2\to B/Q_1\cong \mb{F}_{5}\]
が引き起こされる. また, 次が可換である.
\[
\xymatrix{
A/P_2 \ar[d]\ar[r] & B/Q_2 \ar[d] \\
A/P_1\ar[r] &B/Q_1
}
\]
$5, x'(1-t) \in Q_2$である. もし$x'\in Q_2$なら, $Q_2\supset (5,x')$.
\[B/(5,x') \cong \mb{Z}[X,T]/(5,X,T^2-6) \cong \mb{F}_{5}[T]/(T-1)\times \mb{F}_{5}[T]/(T+1) \]
だから$(5,x')$を含む素イデアルは$(5,x',t+1), (5,x',t-1)$で, $Q_2\neq Q_1$だから$Q_2 = (5,x',t-1)$となるが, これでは$Q_2\subset Q_1$が満たされないので矛盾. \par
$x'\notin Q_2$とする. このとき$1-t\in Q_2$で$t-1 \in Q_2 \subset Q_1$だから$(t+1) - (t-1) = 2\in Q_2$で$Q_2$が$5\mb{Z}$の上にあることに反する. \qed
\end{egans}

\begin{rem}
整拡大を外した場合, 反例がある. $\mb{Z}\subset \mb{Q}$など(左辺と右辺で次元が異なるため).
\end{rem}

\begin{egprob}

\end{egprob}








\subsection{整閉包}
専門問題で整閉包を求める問題は頻出である.
$A\subset B$を環の拡大とするとき, $B$の元で$A$上整なもの全体を$B$における$A$の整閉包という.
\begin{eg}
\ben
\item $\mb{Q}$における$\mb{Z}$の整閉包は$\mb{Z}$である. ($\mb{Z}$が整閉整域であることによる.)
\item $A$を整域とするとき, $K=\mathrm{Frac}A$における$A$の整閉包$C$が考えられる. $A$が整閉整域であるということの定義は, $C=A$であることである. ユークリッド整域, PID, UFDはどれも整閉整域である. たとえばUFD$R$に対して$R[x_1,x_2,\cdots, x_n]$の商体は$R(x_1,\cdots, x_n)$であって, この商体における$R[x_1,\cdots, x_n]$の整閉包は$R[x_1,\cdots, x_n]$自身である.
\item (\tf{整閉整域でないもの}) この例は専門とする人なら抑えたい. $k[T]$の部分環 $A=k[T^2,T^3]$を考えると$A$は整閉整域ではない. \aka{実際, $A$の商体もまた$k(T)$であって, この中の元$T$は$A$上整である.} よって整閉包は$A$より真に大きい. なお, $A = k[X,Y]/(X^2-Y^3)$と表すこともできるが, このイデアルが定める曲線$Y^3=X^2$はカスプ特異点$(0,0)$を持つ. この特異点の存在が$A$の性質に影響していると考えられているそうである.
\een
\end{eg}
商体における整閉包を考えるときに重要なテクニックについてまとめておく.
\tbox{
\begin{point}
\ben
\item UFD上の多項式環は整閉であることは事実として認められるだろう. このため, たとえば$\mb{C}[x^2,xy,y^2]$など, \tf{UFD上の多項式環}の商体における整閉包は必ず多項式環に含まれると言える. これにより, 整閉包を上からおさえることができる.
\item \tf{自己同型}は有用である. 自己同型を用いることで, 商体が有理関数体の代数拡大であるときには トレースやノルムがまた整閉包に属すことを用いるとよい. 超越的
\een
\end{point}
}{}

\begin{eg}
\ben
\item $\mb{C}[x^2,y^2]$の商体における整閉包$R$を求める($R\cong \mb{C}[X,Y]$で整閉なので答えは分かってるが, 遠回りでこれを確かめる).  $R$の元は$\mb{C}(x^2,y^2)$に含まれかつ$\mb{C}[x^2,y^2]$上整なので, とくに$\mb{C}[x,y]$上整である. $\mb{C}[x,y]$は整閉だから,  その元は$\mb{C}[x,y]$に含まれる. よって$R\subset \mb{C}[x,y]\cap \mb{C}(x^2,y^2)$である. さて, $\mb{C}(x^2,y^2)$には$x\mapsto \pm x, y\mapsto \pm y$による自己同型があり, これらの自己同型により固定される$\mb{C}[x,y]$の元は$\mb{C}[x^2,y^2]$に属す元である. よって$R=\mb{C}[x^2,y^2]$である.
\item $\mb{Z}$の$\mb{Q}(\sqrt{2})$における整閉包を求める. $\mb{Q}(\sqrt{2})$は自己同型$\sqrt{2}\mapsto -\sqrt{2}$を持つ. $p+q\sqrt{2}$が整閉なら, $2p, p^2-2q^2$は$\mb{Z}$上整な有理数だから整数である. よって$p\in \frac{1}{2}\mb{Z}$である. もし$p$が整数でなければ$p^2-2q^2$が整数にならないことが簡単に分かる. よって$P\in \mb{Z}$で$2q^2\in \mb{Z}$なので$q$も整数であると分かる.
\item $\mb{C}[x, \sqrt{x^3 + 1}] \cong \mb{C}[X,Y]/(Y^2 - X^3 - 1)$の商体における整閉包を調べる. $y=\sqrt{x^3+1}$とおく. $y$は$\mb{C}[x]$上代数的であり, $\mb{C}(x)$上の共役元は$-y$である. 商体は$\mb{C}(x,y)$だが, $y$が$\mb{C}(x)$上代数的なので$\mb{C}(x,y) = \mb{C}(x)[y]$である. これは$\mb{C}(x)$上の二次拡大体であるから, この中の元は$yf(x) + g(x)$と表せる($f,g\in \mb{C}(x)$). さて, この元が$\mb{C}[x,y]$上整であるとしよう. $\mb{C}(x,y)$には$y\mapsto -y$による自己同型があるから$g(x)-yf(x)$も整である. よって$2g(x)$が$\mb{C}[x,y]$上整であるが, \ao{$\mb{C}[x]\subset \mb{C}[x,y]$は整拡大なので} $2g(x)$は$\mb{C}[x]$上整であり, 整閉性から$g(x) \in \mb{C}[x]$が従う. また, $yf(x)$が整だから$y^2f(x)^2 = (1+x^3)f(x)^2$も整である. $1+x^3$に平方因子が無いことから, $f(x)\in \mb{C}[x]$でなければならない. よって$\mb{C}[x,\sqrt{x^3+1}]$は整閉である.
\een
\end{eg}
\begin{prac}
上の(3)の議論を, 体$k$, 平方因子の持たない$n$変数多項式$F(x_1,\dots, x_n)\in k[x_1,\dots, x_n]$, について$k[x_1,x_2,\dots, x_n, \sqrt{F(x_1,\dots, x_n)}]$の場合に一般化し, これが整閉であることを示せ. $k$を\tf{UFD}にしても議論が回ることを確認せよ. \\
(未検証) UFDでない整閉整域の場合などはどうなるだろうか?
\end{prac}


\begin{prac}
$A\subset B$を整拡大とする. $B$の素イデアル$\maf{q}$を取るとき, $A/\maf{q}\cap A \subset B/\maf{q}$とみなせる. これは整拡大であることを示せ.
\end{prac}

\begin{egprob} (H22数学II,2) $R = \mb{C}[x,y,z]/(y^3-x^2z)$とおく.
\ben
\item $R$は整域であることを示せ.
\item $R$の商体$Q$は$\mb{C}$上純超越拡大であることを示せ.
\item $R$の$Q$における整閉包$\witi{R}$を求めよ.
\een
\end{egprob}
\begin{egans}
$y^3-x^2z$の形を見れば, $f(x,y,z)$を$y$を変数に見て割るのが都合がいい. $x,y,z$の$R$における像も同じ記号で表すと, すべての$R$の元が$y^2f_1(x,z) + yf_2(x,z) + f_3(x,z)$の形に書くことができる. \\
(1): $\phi:R\to \mb{C}[t,s]$を$\phi(x) = t^3, \phi(z) = s^3$, $\phi(y) = t^2s$と定めるとwell-definedである. $R$の元を先のように書き, これが$\ker{\phi}$に含まれるとすると
\[t^4s^2f_1(t^3,s^3) + t^2sf_2(t^3,s^3) + f_3(t^3,s^3)\]
であり, $t,s$の次数を考えれば$f_1,f_2,f_3 = 0$が従うので$\phi$は単射で, $R$は整域の部分環に同型である. \\
(2): (1)より$R\cong \mb{C}[t^3,s^3,t^2s]$であるから, 右辺について調べる. $Q\cong \mb{C}(t^3,s^3,t^2s) = \mb{C}(\frac{t}{s},t^3)$であり, $t^3$は$\mb{C}(\frac{t}{s})$上超越的かつ$\frac{t}{s}$は$\mb{C}(t^3)$上超越的なので, $\set{t^3,\frac{t}{s}}$は$\mb{C}$上代数的独立であり, 純超越基となる.
\begin{recall}
超越拡大$L/K$と$S\subset L$に対し, $S$(の任意の有限部分集合)が代数的独立であり, かつ $L/K(S)$が代数拡大であるとき$S$を超越基底といい, $L=K(S)$のとき純超越基底という.
\end{recall}
(3): $\witi{R}$は$\mb{C}[t^3,s^3,t^2s]$の$\mb{C}(\frac{t}{s},t^3)$における整閉包である. この商体は$\zeta = \zeta_3$として$t\mapsto \zeta t, s\mapsto \zeta s$という自己同型を持つから, $\witi{R}$はこの自己同型で不変な元のみからなる. たとえば$t^3,s^3,t^2s$のみならず, $ts^2$という元も不変であるが, $ts^2$は$X^3 - t^3s^6$の根だから$\witi{R}$に属す. よって$\mb{C}[t^3,t^2s,ts^2,s^3]\subset \witi{R}$が従う. $\witi{R}$の任意の元は$\mb{C}(t,s)$に含まれ, $\mb{C}[t,s]$上整であるから整閉性によって$\witi{R}\subset \mb{C}[t,s]$である必要がある. そして$\mb{C}[t,s]$の元のうち先の自己同型で不変なものは$\disp\oplus_{i+j\equiv 0\pmod{3}} \mb{C}x^iy^j = \mb{C}[t^3,t^2s,ts^2,s^3]$に属すので, $\witi{R}\subset \mb{C}[t^3,t^2s,ts^2,s^3]$が従う.
\end{egans}







\clearpage
\subsection{演習問題}

\subsubsection{Level B}
\begin{prob} (H15数学II,1)
\ben
\item $\mb{Z}[\sqrt{5}]$は整閉でないことを示せ.
\item $p,q$を平方院試を含まない1と異なる奇数で$p\neq q$とするとき, $R=\mb{Z}[\sqrt{p},\sqrt{q}]$は整閉でないことを示せ.
\een
\end{prob}


\begin{prob} (H15数学II,2)
\ben
\item 可換環$R$は整域$A$の部分環とする. 商体の拡大$Q(R)\subset Q(A)$が代数的ならば, $A$の0でないイデアル$I$について$I\cap R \neq 0$が成り立つことを示せ.
\item (1)において, $Q(A)/Q(R)$が代数拡大という条件をはずすと主張が必ずしも成り立たないことを, 例をあげて示せ.
\item 可換環$A$は部分環$R$上の整拡大とする. $P$は$R$の素イデアルとし, $A$の相異なる素イデアル$Q_1,Q_2$が$R\cap Q_1 = R\cap Q_2 = P$を満たすとする. $Q_1,Q_2$には含む含まれるの関係はないことを示せ\footnote{この問題は\cite{atiyah}の5章にある. 「整拡大のときにファイバーの点に埋没はない」というスローガンで覚えている.}.
\een
\end{prob}

\begin{prob} (H31専門1)
$A=\mb{R}[X.Y]/(X^2+Y^2)$とおく.
\ben
\item $A$は整域であることを示せ.
\item $A$の商体を$K$とおき, $A$の$K$における整閉包を$B$とおく. $A$加群としての$B$の生成系を一組与えよ.
\een
\end{prob}

\subsubsection{Level C}
\begin{prob} (R2専門科目) \label{prob:R2senmon3}
整域$A$に対する次の性質を考える.
\lefba{$(\ast)$ $A$の0でない素イデアルのうち極小なものは単項イデアルである.
}
$A$をNoether整域, $x\in A$を0でない$A$の素元とする. このとき, $A[1/x]$が性質を持てば, $A$も性質$(\ast)$を持つことを示せ.
\end{prob}




\subsubsection{Hint・Answer}
{\small
\begin{probans}
(1): $\frac{1+\sqrt{5}}{2}$が入ってないから. \\
(2):  $a=\frac{1+\sqrt{p}+\sqrt{q}+\sqrt{pq}}{2}$を考えよ.
%\[a^2 = \frac{1+p+q+pq}{4} + \frac{q+1}{2}\sqrt{p} + \frac{p+1}{2}\sqrt{q} + \sqrt{pq} \in R \]
\end{probans}

\begin{probans}
(1): $a\in I$をノンゼロとする. $a\in Q(A)$は$Q(R)$上代数的だから, ある$t_1,\dots, t_{n-1}\in Q(R)$が存在して
\[a^n + t_{1}a^{n-1} + \dots + t_{n-1} = 0\]
が成り立つ. $n$が最小であるものと仮定してよく, このとき($a\neq 0$および整域であることを用いて)$t_{n-1}\neq 0$である. $t_i = r_i/s_i$ ($r_i,s_i\in R, s_i\neq 0$)と表し, $s=s_1\dots s_{n-1}$とすると, $st_{n-1} \in R \cap I$が分かる. $s,t_{n-1}\neq 0$なので, 示せた. \\
(2): $R=\mb{C}, A=\mb{C}[T]$とおく. $I=TA$とする. $Q(A) = \mb{C}(T)$, $Q(R) = \mb{C}$なので商体は超越拡大. そして$I\cap R=0$である. \\
(3): $Q_1\subsetneq Q_2$であると仮定して矛盾を導けばよい. $R\to A$は整拡大であるから, $R/P\to A/Q_1$も整拡大である. $Q(A/Q_1)\supset Q(R/P)$が代数拡大であることを証明する. $x\in A/Q_1$は$R/P$上整だから $x/1$は$Q(R/P)$上代数的である. さらに$x\neq 0$であるとき, \ao{$x/1$の$R/P$係数代数関係式を用いることで, $1/x$の$R/P$係数代数関係式を得る.} よって$1/x$も$Q(R/P)$上代数的であるから, 任意の$s/t\in Q(A/Q_1)$が$(s/1)(1/t)$と書けることにより$Q(A/Q_1)\supset Q(R/P)$は代数拡大である. さて, (1)によれば$A/Q_1$のイデアル$Q_2/Q_1$の制限$R/P \cap Q_2/Q_1$は0ではないが, $Q_2\cap R = P$なので$R/P \cap Q_2/Q_1 = 0$であり, 矛盾である. よって包含はない. \qed
\end{probans}

\begin{probans}
(1): $\mb{R}$代数の射$\phi:\mb{R}[X,Y]\to \mb{C}[T]$を$\phi(X) = T$, $\phi(Y) = iT$で定める. $\ker{\phi} \supset (X^2+Y^2)$は明らか. $\mb{R}[X,Y]$の元を$f(X) + Yg(X) + (X^2+Y^2)P(X,Y)$と表す. ($f,g\in \mb{R}[T]$).もしこれが$\ker{\phi}$に属せば$f(T) + iTg(T) = 0$だが$f,g$が実係数なので$f,g=0$となり$\ker{\phi}\subset (X^2+Y^2)$となる. よって$A$は整域の部分環に同型なので整域である.\qed \\
(2): (1)より$A\cong \mb{R}[T,iT] = \mb{R}\oplus \bigoplus_{k\geq 1} \mb{C}T^k$である. $K$は$\mr{Frac}\mb{C}[T] = \mb{C}(T)$に含まれており, $iT/T = i$を含むので$K=\mb{C}(T)$である. $i\in K$は$i^2 + 1 = 0$という$A$係数の整等式を満たすので$i\in B$である. よって$\mb{C}[T]\subset B$である. $B$は$A$上整なので, とくに$\mb{C}[T]$上整である. $B\subset K$だから, $\mb{C}[T]$の整閉性より$B=\mb{C}[T]$である. 明らかに, $B=1A + iA$である.
\end{probans}
%=============================Level C=============================
\begin{probans}
$A[1/x] = A_x$である. $\maf{p}\in \spec{A}$を0でない極小な素イデアルとする. もし$x\in \maf{p}$なら, $(x)\subset \maf{p}$だが, $x$が素元だから$(x)\in \spec{A}$である. よって極小性により$\maf{p}=(x)$で, この場合はよい. 次に$x\notin \maf{p}$とする. このとき$\maf{p}A_{x} \in \spec{A_x}$は極小である. なぜなら, もし$Q \subset \maf{p}A_{x}$なる$Q\in \spec{A}_{x}$があれば, $Q$は$\maf{q}\in \spec{A}$で$x\notin \maf{q}$なるものの拡大$\maf{q}A_x$として得られるから$\maf{q}A_{x} \subset \maf{p}A_{x}$である. これを$A\to A_{x}$で引き戻すと$\maf{q}A_x\cap A = \maf{a} \subset \maf{p}A_{x} \cap A = \maf{p}$なので極小性から$\maf{p}=\maf{q}$であり, $\maf{q}A_x = \maf{p}A_{x}$にもなるから. よって条件より$\maf{p}A_x$は単項イデアルであり, $s/x^n A_{x}$と表せる. $1/x^n$は単元なので$(s/1) A_{x}$とも書ける. \par
\aka{ここで, $s\notin (x)$であると仮定できることを示そう.} もし$s\in (x)$なら$s=xs_1$と表せるが, $s_1/1$と$s/1$は同伴元であるから生成するイデアルは変わらない. もしこれで$s_1\notin (x)$となればこの$s_1$に取りかえる. もし$s_1\in (x)$なら$s_1 = xs_2$とする. この操作は無限回続けられるということはない. 実際, もしこのようにして$s_n = xs_{n+1}$となる$s_n$たちが取れれば, $A$のイデアル増加列
\[(s) \subset (s_1) \subset (s_2) \subset \dots\]
が得られるが, Noether性によりこの列は停留する. したがって$(s_n) = (s_{n+1}) = (xs_{n})$となる$n$があるので, 単元$\epsilon\in A^{\times}$により$s_{n} = \epsilon xs_n$と表せる. $s_n(1-\epsilon x) = 0$で, もし$s_n \neq 0$なら整域であることから$\epsilon x = 1$で$x$が単元かつ素元となるので不適. よって$s_n = 0$であるが, このとき$s=0$であり
, $\maf{p}A_{x} = 0$となる. これは$\maf{p}\neq 0$に反する. したがって$s$を$s_n \notin (x)$となるような$s_n$に取り替えればよい. \par こうして$s\notin (x)$とできる. このとき$\maf{p}=(s)$であることを示そう. \par
まず$s\in \maf{p}$を示す. [証明] $s/1 = r/x^n$となる$r\in \maf{p}$, $n\in \mb{N}$があるので, $x^m(sx^n - r) = 0$なる$m$がある. 右辺は$\maf{p}$に含まれ, $x\notin \maf{p}$なので$sx^n - r\in \maf{p}\implies sx^n \in \maf{p} \implies s\in \maf{p}$. したがって$(s) \subset \maf{p}$が示せた. \par
逆に$r\in \maf{p}$を任意に取る. $r/1\in \maf{p}A_{x} = (s/1)A_{x}$だから$r/1 = (st)/x^n$と表せる. このような$n$で最小のものを取る. $x^{m}(x^n r - st) = 0$である. 整域上であるから, $x\neq 0$より $x^n r = st$である. $n\geq 1$であるなら,  $st \in (x)$かつ$s\notin (x)$かつ$x$が素元であることより$t\in (x)$が従う. よって$t=xt_1$と書け, $x^{n-1}r = st_1$で, $x$の指数が0になるまでこれを続けることができ, 最後に$r = st'$となる$t' \in A$が存在することが示される. よって$r\in (s)$だから$\maf{p}\subset (s)$である. \qed
\end{probans}

}



\newpage





















\section{加群論}
\subsection{加群の種類}
\tbox{
\begin{defn}$A$は環, $M$を$A$加群とする.
\ben
\item 自由加群
\item 有限生成加群
\item 平坦加群
\item 射影加群
\item Noether加群
\item ($\spadesuit$ ) 入射加群
\item ($\spadesuit$) Artin加群
\een
\end{defn}
}{title=加群の種類}
自由加群が最もベクトル空間に近く, かなり性質がよい. 平坦加群や射影加群も, 単因子論などで用いられる.

\tbox{
\begin{eqnarray*}
\text{自由}\subset \text{平坦} \subset \text{射影}
\end{eqnarray*}
}{title=加群のクラス}

射影加群や平坦加群については, ホモロジー代数の節で少し取り扱う.

\begin{prop} $M$は$A$加群とする.
\ben
\item
\item
\item $A$がNoether環で$M$が有限生成であるとき, $M$はNoether加群である.
\een
\end{prop}






\begin{egprob}
何の問題か忘れた
\end{egprob}

\begin{egans}
(1): $x=b/a + \mb{Z}$とおく。ただし$a,b\in \mb{Z}$で$a\neq 0$である。$ax = a+\mb{Z} = \mb{Z}$なので$ax = 0$である。\qed\\
(2): $x=s/t + \mb{Z}$とおく。ただし$s,t\in \mb{Z}$, $t\neq 0$である。$y=s/(ta ) + \mb{Z}$とおけば$ay = x$\qed \\
(3): 自由加群ではない。背理法で示す。もし$\mb{Q}/\mb{Z}$に$\mb{Z}$基底$\{ e_i\}_{i\in I}$が存在したとすれば, それらは非自明な線型関係式を持たないが, (1)から$a_i\in\mb{Z}\setminus{\{ 0 \}}$であって $a_i e_i = 0$となるものが存在する。この式自体は$a_i\neq 0$によって非自明な線型関係式であるから矛盾である。\qed

(4):有限生成ではない。背理法で示す。$\mb{Q}/\mb{Z}$が$x_1,\cdots,x_n\in \mb{Q}/\mb{Z}$で$\mb{Z}$上生成されたと仮定する。このとき, 任意の素数$p$に対して, ある$a_{i,p}\in \mb{Z}$ ($i=1,\cdots, n$)が存在して
\[\dfrac{1}{p} + \mb{Z} = \sum_{i=1}^{n} a_{i,p}x_i\]
と書ける。ここで, (1)からある$b_i\in\mb{Z}\setminus{\{ 0 \}}$が存在して$b_ix_i \in \mb{Z}$となるものが存在するので, $b=b_1\cdots b_n$として
\[\dfrac{b}{p} + \mb{Z} = \sum_{i=1}^{n}a_{i,p}bx_i = 0\]
である。すなわち, 任意の素数$p$に対して$\dfrac{b}{p} \in \mb{Z}$である。しかし$b$の素因数でない$p$を取ることでこのようなことは起こりえない。よって矛盾である。\qed
\end{egans}

$A$加群の準同型$\phi:M\to N$に対し, $\ker{\phi}, \im{\phi}, \coker{\phi}$というものが定義される.


\begin{egprob} (\cite{aoyukie}, 演習2.5)
$A= \mr{diag}(1,3,6,0)$とし, $T_A:\mb{Z}^4\to \mb{Z}^4$を$A$により定義される加群の準同型とする. このとき, $\ker{T_A}$, $\im{T_A}$, $\coker{T_A}$を求めよ.
\end{egprob}
\begin{egans}
\[\im{T_A} = \mb{Z}\oplus (3\mb{Z})\oplus 6\mb{Z} \oplus 0\]
\[\ker{T_A} = 0\oplus 0\oplus 0\oplus \mb{Z} \]
\[\coker{T_A} = 0\oplus \mb{Z}/3\mb{Z} \oplus \mb{Z}/6\mb{Z} \oplus \mb{Z}\]
\end{egans}



有限生成加群といえば, \tf{中山の補題}は大切である.

\tbox{
\begin{prop}
$M$を有限生成$A$加群, $I$を$A$のJacobson根基に含まれるイデアルとする. $IM = M$であるとする. このとき, $M=0$である.
\end{prop}
}{title=中山の補題}

\begin{prac} (仮想面接問題！これは聞かれそう)
有限生成を外した場合, 中山の補題の反例は?
\end{prac}













\subsection{単因子論とスミス標準形}
次の二つはとても重要である.
\begin{theo}
$R$がPID, $M$が有限生成$R$加群なら, 自然数$r$と$e_1,\cdots, e_t\in R$で次を満たすものが存在する。
\ben
\item $M\cong R^{r} \oplus R/(e_1) \oplus R/(e_2) \oplus \cdots \oplus R/(e_t)$
\item $e_i | e_{i+1}$ ($i=1,2,\cdots, t-1$)
\een
また上の条件が満たされるなら, $r$は一意であり, イデアル$(e_1),\cdots, (e_t)$も一意である。
\end{theo}

\begin{cor}
有限生成$R$加群$M$($R$はPID)に対し, $M$が平坦であることとTorsion Freeであることは同値である.
\end{cor}
\prf Torison Freeなら$M \cong R^r$なので, 自由加群だから平坦になる. 逆に$M$が平坦であると仮定しよう. $a\neq 0 \in R$に対して$\phi_a: R\to R$を$\phi_{a}(r) = ar$と定めると, \tf{$R$が整域だから}これは単射である. よって, 平坦性により$\tens_{R}M$は左完全だから, 単射性が保存される. つまり, 次も単射である.
\[\phi_{a} \tens \mr{id}_{M} : R\tens_{R} M \to R\tens_{R} M\]
$\Phi: R\tens_{R} M \cong M$の同型写像と合成すると, この写像は
\[m_{a}: M\to M,\quad m(x) = ax\]
という$a$倍写像になる. $\phi_{a}\tens \mr{id}_{M}$と$m_a$は$\Phi$で移り合うので(図式参照), $m_a$も単射である. これが任意の$a\in R$で成り立つということは, $M$のねじれ元が$0$のみであるということに他ならない. \qed

\begin{eg}
$\mb{Z}$加群$\mb{K} = \mb{Q}, \mb{R}, \mb{C}$などはTorison Freeなので平坦である. よって関手$\cdot  \tens_{\mb{Z}} \mb{K}$右完全であるばかりでなく, 完全列を保つ. とくに, \tf{テンソルと剰余の交換が可能である。}(これは一般には成り立たないはず.) より一般には, $F$が平坦$R$加群, 部分加群$N\subset M$とするとき,
\[(M\tens_{R} F)/(N\tens_{R} F) \cong (M/N)\tens_{R} F \]
である(適当な完全列を取って示してみよ).
\end{eg}

\begin{prac}
\ben
\item 自然数$n$に対して$(\mb{Z}/n\mb{Z})\tens_{\mb{Z}} \mb{K}  = 0$を示せ. (\tf{Hint}: 剰余とテンソルの交換を用いよ.)
\item $M$をrank $n$の有限生成$\mb{Z}$加群とする:
\[M\cong \mb{Z}^{n} \oplus T\]
$n=\dim_{\mb{K}} (M\tens_{\mb{Z}} \mb{K})$を示せ.
\een
\end{prac}

\begin{rem}
とくに$R=\mb{Z}$の場合は有限生成Abel群の構造定理と呼ばれる定理になる.
\end{rem}

\tbox{
\begin{theo}
$R$をPIDとする(たとえば$R=\mb{Z}$). 任意の$A\in \mr{M}_{n}(R)$に対して, ある$P,Q\in GL_{n}(R)$が存在して,
\[PAQ = \bep
e_1 & & &  \\
& \ddots & &&\\
&&e_r &  &\\
&&&0&\\
&&&&\ddots
\enp\]
\end{theo}
とできる. ここで$e_1,\cdots, e_r\neq 0$かつ$e_1 | e_2 | \cdots | e_t$であり, このような$e_i$は同伴を除いて一意に定まる. これら$e_i$を行列$A$の単因子といい, 右辺の行列はスミス標準形, あるいは単因子標準形と呼ばれる.
}{title=スミス標準形の存在性}

\begin{rem}
$P,Q$は基本行列の積として取れる. つまり, 行列の基本変形でスミス標準形に持って行くこともできる. しかし線型代数のときのような「第$k$列を$d$ ($d\neq 0,\pm 1$)で割る」操作に対応する基本行列を用いることはできない. $GL_{n}(\mb{Z})$の元は行列式が$\pm 1$でなければならないからだ. このように, 線型代数の場合とは使用できる基本変形が大きく異なることに注意しなければならない. 使用できる操作をまとめると
\ben
\item 第$i$行に第$j$行の\tf{整数倍}を足す操作
\item 第$i$行を$\pm 1$倍する操作
\item 行同士を交換する操作
\een
および上の二つで「行」を「列」に変えた操作である.
\end{rem}


基本行列の積で表せるとは言え, 基本変形だけでは少し難しいこともある. 以下は整数行列の変形法をまとめたものである(\cite{aoyukie} 139p参照).
\tbox{
\ben
\item $A$が零なら何もしない. $A$が零でないなら, 0でない成分の中で\tf{絶対値が最小のもの}を行や列を交換して$(1,1)$成分に移す.
\item もし第1行や第1列で$(1,1)$成分で割り切れないものがあれば,
\[ P =
\bep
x & y & \\
-q & p & \\
& & I_{m-2}
\enp
\]
のような行列を使う. ただし$px+qy=1$である\footnote{$p,q$が与えられており, $x,y$は$px+qy=1$の一つの解として求める.}. 0でない成分のうち, 絶対値がより小さいものを作る. (1)を繰り返し, 第1行と第1列の$(1,1)$成分以外が0になるまで続ける. $\pm 1$をかけ, $(1,1)$成分が正になるようにする.
\item ほかに$(1,1)$成分で割り切れないものがあれば, その成分を含む列を第1列に加え, (1),(2)を行う. これを繰り返せば, (2)の状態でさらにすべての成分が$(1,1)$成分で割り切れるように変形できる.
\item (1)-(3)のステップを第1行と第1列を除いた行列に繰り返す.

\een
}{title=整数行列の変形方法}

\begin{eg} (\cite{aoyukie} 例2.12.7)
$R=\mb{Z}$, $A=\bep 24 & 32 \\ 10 & 14 \enp$とするとき, $\coker{T}_{A}$を巡回加群の直和に表す. $A$の成分は偶数で,第1列は$2[12,5]$で, $12$と$5$は互いに素. $(-2)\cdot 12 + 5\cdot 5 = 1$となるから, $(p,q,x,y) = (12,5,-2,5)$により
\[P = \bep -2 & 5 \\ -5 & 12 \enp \in SL_{2}(\mb{Z})\]
を考えると,
\[2P \bep 12 & 16 \\ 5 & 7 \enp = 2 \bep 1 & 3 \\ 0 & 4 \enp\to \bep 2 & 0 \\ 0 & 8 \enp\]
となる. ただし最後は, 第1列の3倍を第2列から引いた. 一般に$GL_{n}(\mb{Z})$の元が左や右からかかっても, ker,im,cokerは同型を除いて不変なので\footnote{左からかかっても不変であるというのは, $T_P$が$im{T_A}$を$\im{T_{PA}}$に移すことから考えるとよい. 詳しくは\cite{aoyukie}命題2.5.7},
\[\coker{T_A} \cong \coker{T_{PA}} = \mb{Z}/2\mb{Z} \oplus \mb{Z}/8\mb{Z}\]
\end{eg}

\begin{egprob}
$A$が次の整数行列であるとき, $\coker{T_A}$を求めよ.
\begin{center}
(1):$\bep 6&12 \\ 4 & 6 \enp $\quad (2):$\bep 3 & 7 \\ 10 & 22\\ 9&21 \enp$\\
(3): $\bep 14 & 11 \\ 6 & 5\\ 2 & 4 \enp$\quad (4): $\bep 5 & 2&5\\6&3&10\\ 10&4&16\\ 11&5&15 \enp$
\end{center}
\end{egprob}
\begin{egans}
(1): 以下のように変形できる.
\begin{eqnarray*}
\bep 6 & 12 \\ 4 & 6 \enp \to \bep 4 & 6 \\ -2 & 0 \enp \to \bep 0 & 6 \\ 2 & 0 \enp
\end{eqnarray*}
よって, \bolm{\coker{T_A} \cong \mb{Z}/2\mb{Z} \oplus \mb{Z}/6\mb{Z}}\\
(2):
\begin{eqnarray*}
\bep 3&7\\10&22\\9&21 \enp \to \bep 3&7\\1&1\\0&0 \enp \to \bep 0&4\\1&1\\0&0 \enp \to \bep 1&1\\0&4\\0&0 \enp \to \bep 1 & 0 \\ 0&4\\0&0 \enp
\end{eqnarray*}
より, \bolm{\coker{T_A}\cong 0 \oplus \mb{Z}/4\mb{Z} \oplus \mb{Z}}である. \\
(3):
\begin{eqnarray*}
\bep 14&11\\6&5\\2&4 \enp \to \bep 0&-17\\0&-7\\2&4 \enp \to \bep 0&3\\0&7\\2&4 \enp \to \bep 2&4\\0&3\\0&1 \enp \to \bep 2&0\\0&0\\0&1 \enp
\end{eqnarray*}
より, $\coker{T_A} \cong \mb{Z}/2\mb{Z} \oplus \mb{Z}\oplus 0$である.

(4): 成分の中で一番小さいのが2なので, それを(1,1)成分に持ってくる.

\begin{eqnarray*}
\bep 2&5&5\\3&6&10\\4&10&16\\5&11&15 \enp \to \bep 2&5&5\\1&1&5\\ 0&0&6\\1&1&5 \enp \to \bep 0&3&-5\\ 1&1&5\\ 0&0&6\\ 0&0&0 \enp\\
\to \bep1&1&5\\0&3&-5\\ 0&0&6\\ 0&0&0 \enp \to \bep 1&0&0\\0&3&-5\\0&0&6\\0&0&0 \enp
\end{eqnarray*}
2,3列, 2,3行の小行列を見ると
\begin{eqnarray*}
\bep 3&-5\\0&6 \enp \to \bep 3&1\\0&6 \enp
\end{eqnarray*}
ここからは右上の1を利用します. 「1」がある列または行に, その「1」以外の非零成分がなければ, その行または列を使って簡約できます.
\[\bep 3&1\\0&6 \enp \to \bep 3&1\\-18&0 \enp \to \bep 0 & 1\\ -18&0 \enp\]
よっ, 元の行列に戻れば\bolm{\coker{T_A} \cong  \mathbb{Z}/18\mathbb{Z} \oplus \mathbb{Z} }となる.
\end{egans}


\begin{egprob}
アーベル群$G=\mb{Z}/2\mb{Z}\times \mb{Z}/4\mb{Z}\times \mb{Z}/6\mb{Z}$の元$a=(1,3,2)$で生成される部分群を$H$とする。剰余群$G/H$を巡回群の直積で表せ。
\end{egprob}

\begin{egprob}

\end{egprob}























\newpage
\subsection{テンソル}

\begin{prop} $M$を$A$加群とするとき,
テンソル関手$\bullet \tens_{A} M$は右完全である.
\end{prop}
\prf
$0\to M_1\to M_2\to M_3 \to 0$を完全とする. $f:M_1\tens_{A} M \to M_2\tens_{A} M$, $g:M_2\tens_{A} M \to M_{3}\tens_{A} M$とする. 一番難しいのは$\ker{g} = \im{f}$の部分である. これは$\witi{g}: (M_2\tens_{A} M)/\im{f} \to M_3\tens_{A} M$が同型であることを示せばよい. そのために, \bolm{\witi{g}の逆写像を作る.} まず, 直積からの射$h:M_3\times M \to (M_2\tens_{A} M)/\im{f}$を次で定める.
\lefba{
$(x,z)\in M_3\times M$に対し, 全射$M_2\to M_3$による$x$のリフトを任意に一つとり, それを$y$とする. これを用いて
\[h(x,z) = y\tens z \bmod{\im{f}} \]
と定める.
}
これのwell-definedを示そう. まず$y_1,y_2$がともに$x$のリフトであるとせよ. このとき$y_1\tens z - y_2\tens z = (y_1-y_2)\tens z$である. $y_1-y_2$は$0\in M_3$のリフトである. つまり$M_2\to M_3$の核に属すので, $p:M_1\to M_2$として, $p(t) = y_1-y_2$と表せる. よって
\[p(t)\tens z = (p\tens \mr{id}_{M})(t\tens z) = f(t\tens z)\in \im{f}\]
なので, well-definedである. また, 明らかに$z$に関しては$A$線形であるし, $x$に関しても$A$線形であることは容易に示せる. したがって$h$は$\witi{h}:M_3\tens_{A} M \to (M_2\tens_{A} M)/\im{f}$を引き起こし, これが$\witi{g}$の逆写像であることが容易に分かる($g$でリフトを選択して定義しているので).\qed

\begin{egprob}(仮想面接問題)
テンソル関手が左完全にならない例をあげてください.
\end{egprob}
\begin{egans}
2倍写像$\mb{Z}\to \mb{Z}$は単射だが$\mb{Z}/2\mb{Z}$にテンソルしたものは零写像なので単射でない. すなわち, $\cyc{2}$は$\mb{Z}$加群として\tf{平坦}ではない.
\end{egans}
左完全は一般に保存されないが, 保存されるものを\tf{平坦加群}という.


\subsubsection*{補足-ねじれなしと平坦-}
{\small
単因子論を用いてPID上の有限生成加群について
\[ねじれなし\iff 平坦\]
を証明した. これは有限生成を外しても成り立つことである. ここではその証明をまとめておく.
\begin{theo} $A$を環とする.
\ben
\item $A$加群$M$が平坦であるための必要十分条件は, 任意のイデアル$I\subset A$に対して$I\tens_{A} M \to IM$が同型であることである.
\item $A$がPIDのとき, $A$加群$M$が平坦であることと$M$がねじれなしであることは同値である.
\een
\end{theo}
\prf
(1): 必要性は明らか. $I\tens_{A}M\to IM$が同型であると仮定する. $N'\to N$を任意の単射とする($N'\subset N$とみなす). この$N$に関して場合分けを行う. 自由加群からの全射準同型$p:F\to N$が存在する.  次の図式を考える.
\[
\xymatrix{
\ker{p}\tens_{A}M \ar@{=}[d] \ar[r]& p^{-1}(N')\tens_{A}M \ar[d] \ar[r] & N'\tens_{A} M\ar[d] \ar[r] & 0\\
\ker{p}\tens_{A} M \ar[r]&F\tens_{A}M  \ar[r]& N\tens_{A}M \ar[r]& 0
}
\]
示したいことは$N'\tens_{A} N' \to N\tens_{A}M$の単射性であるが, このためには$p^{-1}(N')\tens_{A} N' \to F\tens_{A}N$が単射であることがいえればよい(diagram chasing またはFive Lemma). \tf{したがって主張は$N$が自由加群の場合に帰着する. } まず$N$を有限ランク$n$の自由加群とし, $n$に関する帰納法で有限ランクの場合を証明する. $n=1$ならば, $N\cong A$で$N'$は$A$のイデアルに同一視されるので仮定より明らか(即ち, この仮定は帰納法の最初のステップのためにある). $n>1$とし, ランク$n-1$以下の自由加群$N$で主張が正しいとせよ. $N_1,N_2\neq N$は $N = N_1\oplus N_2$なる直和因子とする. この分解を$N'$に落とす:
\[N_1' = N_1\cap N',\quad N_2':= (N'+N_1)/N_1 \]
そして次の図式を考える(横に完全である).
\[
\xymatrix{
N_1'\tens_{A}M \ar[r]\ar[d]^{\beta} &N'\tens_{A}M\ar[d]\ar[r]& N_2'\tens_{A}M\ar[d]^{\gamma} \\
N_1\tens_{A}M \ar[r]^{\alpha} & N\tens_{A}M \ar[r] & N_2\tens_{A}M
}
\]
ここで$\beta,\gamma$は帰納法の仮定より単射であり, $\alpha$は入射$N_1\to N$からテンソル積によって得られる写像なので単射である. ここからdiagram chasingにより$N'\tens_{A}M \to N\tens_{A}M$も単射であることが分かる.
\begin{prac}
このdiagram chasingをせよ.
\end{prac}
よって有限ランクの場合は証明できた. あとは$N$が自由で任意ランクのときを示せばよい. $N_0$を$M$の有限ランクの直和因子となる自由加群とする. 有限ランクの場合により, $(N'\cap N_0)\tens_{A}M \to N_0 \tens_{A}M$は単射である.
入射$j:N_0\to M$についてはテンソル後も単射であるから, $N_0\tens_{A}M \to N\tens_{A}M$も単射で, したがって合成\bolm{(N'\cap N_0)\tens_{A}M \to N\tens_{A}M}も単射であり, この写像は$N'\cap N_0 \to N$から引き起こされるものに一致する. 実はこれで$N'\tens_{A}M \to N\tens_{A}M$の単射性が得られるのである.

 [証明] $x\in N'\tens_{A}M$がこの写像で送って$0$であるとする. 上手く$N_0$を選べば, \ao{$(N'\cap N_0)\tens_{A}M \to N'\tens_{A}M$の像に$x$が入るようにできる(そのリフトの一つを$y$とおく). }この写像に$N'\tens_{A} M \to N\tens_{A}M$を合成して, \bolm{(N'\cap N_0 )\tens_{A}M \to N'\tens_{A}M \to  N\tens_{A}M}を得るが, これは単射であった(図式の可換性に注意せよ). そして, $y\mapsto x\mapsto 0$と移るので, もともと$y=0$であり$x=0$となる. これで示せた.
 \[
\xymatrix@C=80pt{
(N'\cap N_0)\tens_{A}M \ar[r]^{有限ランクの場合より単射}\ar[d] & \aka{\bolm{N_0\tens_{A}M}}\ar[d]^{入射から来るので単射} \\
N'\tens_{A}M \ar[r]^{Question.単射か?}& N\tens_{A}M
}
\]
(2): $I=aA$を$a\neq 0$とする. $t_a:A\to A$, $u_a: M\to M$を$a$倍写像とする. $t_a:A\to I$にすると(整域だから)同型である. $f:I\tens_{A}M \to IM$を$f(r\tens m) = rm$として, 次は可換.
\[
\xymatrix@C=80pt{
\ar[d]^{u_a} M\cong A\tens_{A}M \ar[r]^{t_a\tens\mr{Id}_M: ここが同型}  & I\tens_{A}M\ar[ld]_{f} \\
IM
}
\]
これより$f$が同型であることと$u_a$が同型であることは同値. $u_a$はもともと全射なのでこれは$u_a$が単射であることに同値で, それは$a\neq 0$倍して消える$M$の元が存在しないということであるから, $M$がねじれなしということに同値である.
\qed





}





\newpage

\subsection{($\spadesuit\spadesuit$)ホモロジー代数}
ホモロジー代数を使う人のために簡単な復習を行う. 以下, $R$は可換環, $A$は$R$代数とし, 加群は$A$上で考える.
\begin{defn}
\ben
\item 関手$\mr{Hom}_{A}(P,\bullet):A-\ken{Mod} \to A-\ken{Mod}$が\tf{右完全}であるような $A$加群$P$を射影的$A$加群という.
\item 関手$\mr{Hom}_{A}(\bullet, I)$が左完全であるような$A$加群$I$を単射的$A$加群(or 入射加群)という.
\een
\end{defn}


\begin{prop} (射影加群について)
\ben
\item 自由加群は射影加群である.
\item 自由加群のある直和因子になる加群は射影加群である. 逆に, 任意の射影加群$P$にはある$N$が存在して$P\oplus N $が自由加群となる.
\een
\end{prop}
\prf (1): $F$を自由加群, $\pi:M\to N$を全射, $\phi:F\to N$を与える. $F=\oplus_{i} Ax_i$とし, $\phi(x_i) = \pi(m_i)$とする. このとき$\psi:F\to M$を$\psi(\sum a_ix_i) = \sum a_im_i$で定めると$\pi\psi = \phi$.

(2): $P$が, ある$Q$が存在して$P\oplus Q = A^{\oplus I}$となるとする. $\pi:M\to N$を全射, $\phi:P\to N$を射とする. $\psi:P\to M$を与えたい. これは次のようにする: $\phi:P\to N$と零写像$Q\to N$を束ねることで$\phi_1: P\oplus Q = A^{(I)}\to N$を得る. $A^{(I)}$は自由加群だから射影的. よって, $\psi_1: P\oplus Q = A \to M$に持ち上がる. $s:P\to P\oplus Q$と$\psi_1$を合成して$\psi:P\to N$を得るが, これで$\pi\circ \psi = \phi$となる. \par
逆に$P$を射影加群とせよ. 自由加群$F$が存在して$F\to P$は全射. $N = \ker{(F\to P)}$とすると$0\to N \to F \to P\to 0$は完全で, $P$が射影的なので分裂する(\tf{分裂補題}). よって$F\cong N\oplus P$.\qed

\qed

\begin{recall}
自由$\implies$射影$\implies$平坦.
\end{recall}
そこで, 聞かれ方がいろいろあるだろう.
\tbox{
\begin{egprob}
\ben
\item 射影加群だが自由加群でない例.
\item 平坦加群だが射影加群でない例.
\item 平坦でない加群.
\een
\end{egprob}
}{}
\begin{egans}
(1): \ao{$A=\mb{Z}/6\mb{Z}$とする.} $A$加群$P=\mb{Z}/2\mb{Z}$を考える. もし$P$が自由加群であるなら, $P\cong A^{\oplus I}$の形だが, これは\tf{位数の都合的に不可能である.} (2が6のべき乗にならないから) 一方で, $A$自由加群$A$は中国剰余定理から$\mb{Z}/2\mb{Z}\oplus \mb{Z}/3\mb{Z}$に同型なので, $P$は\tf{自由加群の直和因子}である. よって射影$A$加群である. \\
(2): \ao{$\mb{Q}$は$\mb{Z}$加群として射影的ではない.} 実際, 平坦なのはねじれがないから. 射影的であるとせよ. 全射$\pi:\mb{Z}^{(I)}\to \mb{Q}$を取る. このとき, $N = \ker{\pi}$とすると$0\to N \to \mb{Z}^{(I)} \to \mb{Q} \to 0$は\ao{$\mb{Q}$が射影的であるという仮定により分裂する}ので$\mb{Z}^{(I)}\cong N\oplus \mb{Q}$となる. これにより$\mb{Q}\to N\oplus \mb{Q} \to \mb{Z}^{(I)}$という単射$f$を得るが, $f(1) = (n_i)_{i\in I}$だとすると$f(1/k) = (n_i/k)_{i\in I}$がすべての$k\in \mb{Z}$で成り立ち, $n_i = 0$でなければならず, $f\equiv 0$なので単射でなくなる. よって射影的ではない. \\
(3):$\mb{Z}$加群$\cyc{2}$. ねじれがあるので平坦でない(単因子論).
\begin{comment}
\ao{$\mb{Z}$加群$\mb{Z}/2\mb{Z}$を考える.} この加群は, $x\bmod{6}$, $y\bmod{2}$に対して$xy\bmod{2}$を与えるという作用で構造を入れる. 目標は, $\pi:M\to N$という全射と, $\phi:\mb{Z}/2\mb{Z}\to N$という射を与え, これが$\psi:\mb{Z}/2\mb{Z}\to M$に持ち上がらないようなものを見つけることである. たとえば$M=\mb{Z}/4\mb{Z}\to \mb{Z}/2\mb{Z} = N$という自然な全射と, $\phi:\mb{Z}/2\mb{Z}\to \mb{Z}/2\mb{Z}$を恒等写像とする. もしこれが$\mb{Z}$加群の射$\psi:\mb{Z}/2\mb{Z}\to \mb{Z}/4\mb{Z}$に持ち上がると, $2\psi(1) = \psi(2) = \psi(0)=0$なので$\psi(1) = 0,2$でなければならず, $\pi\psi(1) = 0 \neq  1 = \phi(1)$だからおかしい. \\
(4):  \qed
\end{comment}
\end{egans}


次は覚えておいてもいいくらいの問題かもしれない.
\tbox{
\begin{egprob} (\cite{kiyukie}問題6.3.1(1))
$A$が可換な整域で, $M$がねじれを持つ$A$加群なら, $M$は射影的加群ではないことを証明せよ.
\end{egprob}
}{}
\begin{egans}
ある零でない$a\in A$により$\phi: M\xrightarrow{\times a}M$は単射ではない. $x\in \ker{\phi}\setminus{\set{0}}$とする. 自由加群からの全射$A^{(I)}\to M$が存在する. この全射と$\mr{id}_{M}$について,  仮に$\psi:M\to A^{(I)}$に持ち上がったとしよう. $ax = 0_M$であるから$\psi(ax) = 0 = a\psi(x)$となる. $\psi(x) = (a_i)_{i\in I}$と書くと, $(aa_i)_{i\in I} = 0$だから$aa_i = 0$がすべての$i$で成り立つ. $A$は整域だから$a_i = 0$である. よって$\psi(x) = 0$を得るので$\pi\psi(x) = \pi (0) = 0_M = \mr{id}_M(x)= x$となり$x\neq 0$に矛盾する. よって$M$は射影的である. \qed

\end{egans}


\tbox{
\begin{prop} (\cite{kiyukie}黄雪江, 定理6.3.11)
$A$は$R$代数, $M$を右$A$加群, $N$を左$A$加群とする. $(P_n,d_{P,n},\epsilon_P)$, $(Q_n,d_{Q,n},\epsilon_Q)$をそれぞれ$M,N$の射影的分解
\[\to P_n\tens_{A} N \xrightarrow{d_{P,n}\tens 1} P_{n-1}\tens_A N \to \dots \to P_0\tens_{A}N \xrightarrow{\epsilon_P} 0\]
\[\to M\tens_{A} Q_n \xrightarrow{1\tens d_{Q,n}} M\tens_A Q_{n-1} \to \dots \to M\tens_{A} Q_0\xrightarrow{\epsilon_Q} 0\]
を考える. このとき,
\ben
\item $R$加群$\mr{H}_n(P_{\ast}\tens_{A} N)$, $\mr{H}_n (M\tens_A Q_{\ast})$は射影的分解$(P_n,d_{P,n},\epsilon_P)$, $(Q_n,d_{Q,n},\epsilon_Q)$の取り方によらず定まる.
\item 右$A$加群$M$と左$A$加群$N$に対し
\[\bolm{\mr{H}_0(P_{\ast}\tens_{A} N) \cong \mr{H}_0(M\tens_{A} Q_{\ast}) \cong M\tens_{A} N}\]
\item $k$加群として\bolm{\mr{H}_n(P_{\ast}\tens_A N) \cong \mr{H}_{n}(M\tens_A Q_{\ast})}.
\een
\end{prop}
}{title=トーションの定義}
\begin{defn}
上の命題の(3)で等しい$k$加群のことを$\mr{Tor}^{A}_{n}(M,N)$と定義し, \tf{トーション}という. (1)より, トーションは$M,N$の片方の射影的分解を任意に一つ選び, もう片方のテンソルにより得られる複体のホモロジーとして得られる.
\end{defn}
双対的にイクステンション$\mr{Ext}^{n}_{A}(M,N)$が定義されるが, これは第一成分と第二成分で扱い方が微妙に異なる. 第一成分$M$では反変的であり, $M$の射影分解を用いるが, 第二成分$N$で共変的であり, \tf{入射分解}を用いる. 定義を覚えておくときのポイントは, 分解でHomを取ったときにそれが\tf{コチェイン複体}であるようにすることである
.

\tbox{
\begin{prop} (\cite{kiyukie} 定理6.3.21)
$(P_n, d_{P,n}, \epsilon_P)$を$M$の射影的分解, $(I_n,d_{I,n},i_I)$を$N$の入射的分解とする., このとき
{\small
\bealn{
0\to \mr{Hom}_{A}(P_0,N) \to \mr{Hom}_{A}(P_1,N) \to \mr{Hom}_{A}(P_2,N)\to \dots \\
0\to \mr{Hom}_{A}(M,I_0) \to \mr{Hom}_{A}(M,I_1) \to \mr{Hom}_{A}(M,I_2)\to \dots
}
}
はコチェイン複体である. この状況で
\ben
\item $R$加群$\mr{H}^n(\mr{Hom}_{A}(P_{\ast},N))$, $\mr{H}^n(\mr{Hom}_{A}(M,I_{\ast}))$は分解の取り方によらず定まる.
\item \bolm{H^0(\mr{Hom}_{A}(P_{\ast},N)) \cong \mr{H}^0(\mr{Hom}_{A}(M,I_{\ast})) \cong \mr{Hom}_{A}(M,N)}
\item $R$加群として\bolm{\mr{H}^n(\mr{Hom}_{A}(P_{\ast},N))\cong \mr{H}^n(\mr{Hom}_{A}(M,I_{\ast}))}. この等しい$R$加群のことを\bolm{\mr{Ext}^{n}_{A}(M,N)}と表し, \tf{イクステンション}という.
\een
\end{prop}
}{title=イクステンションの定義}
上により, $\mr{Ext}^{n}_{A}(M,N)$を計算するだけなら, $M$の射影分解を一つ選択するか, $N$の入射分解をひとつ選択して, Homを取ってコホモロジーを計算するだけでよい. Torはテンソル複体の完全度をはかり, ExtはHom複体の完全度を測る. したがって, この群が0になるというのは嬉しいことである.  これに関連して次は覚えておくべきだろう.
\tbox{
\begin{egprob}
\ben
\item $M$が右$A$加群, $N$が左$A$加群とするとき, \ao{$M,N$のどちらかが射影的}なら, すべての$n>0$で$\mr{Tor}_{n}^{A}(M,N) = 0$
\item $M,N$を加群とするとき, $M$が射影的か, $N$が入射的であるなら, すべての$n>0$に対し$\mr{Ext}^{n}_{A}(M,N) = 0$である.
\een
\end{egprob}
}{}
\begin{egans}
(1): $M$が射影的なら, $M$の射影分解として
\[\dots\to 0\to 0\to P_0 = M \xrightarrow{\epsilon_P} \to 0\]
が取れ, これにテンソルをすれば1次以上が0のチェイン複体ができるので, それのホモロジーは0. \\
(2):同様である. $M$側では射影分解を使って定義し, $N$側では入射分解を用いていたことに注意.
\end{egans}


\begin{egprob}
$A$を環, $x$を零因子でも単元でもない元とする. このとき$M=A/xA$は射影加群ではないことを示せ.
\end{egprob}
\begin{egans}
自然な全射$\pi:A\to A/xA$を与え, $\phi:A/xA\to A/xA$を恒等写像とする. これが$\psi:A/xA\to A$に持ち上がったとする. $\psi(1\bmod{xA})\in A$は,
\[x\psi(1+xA) = \psi(xA) = 0\]
を満たすので$x$が零因子ではないことより$\psi(1+xA) = 0$である. よって$\pi(1 + xA) = \pi\psi(1+xA) = xA$だから$1+xA = xA$となる. よって$xA=A$だが$x$が単元でないからこれはおかしい.
\end{egans}

\tbox{
\begin{egprob}  $n$は0以上の整数とする. 次を求めよ.
\ben
\item $\mr{Tor}_{n}^{\mb{Z}}(\mb{Z},N)$, $\mr{Tor}_{n}^{\mb{Z}}(M,\mb{Z})$
\item $\mr{Tor}_{n}^{\mb{Z}}(\cyc{a},\cyc{a})  $
\item 互いに素な$a,b>0$に対し, $\mr{Tor}_{n}^{\mb{Z}}(\cyc{a},\cyc{b})$
\een
\end{egprob}
}{title=トーション計算問題}

\tbox{
\begin{egprob} $n$は0以上の整数とする. 次を求めよ. \ao{必要ならば, $\mb{Q}/\mb{Z}$が入射加群であることを用いて良い. }
\ben
\item $\mr{Ext}^{n}_{\mb{Z}}(\cyc{2},\cyc{2})$
\een
\end{egprob}
}{title=イクステンション計算問題}
\begin{egans}
(1):次は入射的分解(番号が上がる分解である).
\[0\to \dfrac{1}{2}\mb{Z}/\mb{Z} \to \mb{Q}/\mb{Z} \xrightarrow{\times 2} \mb{Q}/\mb{Z}\to 0\to 0\to 0\to \dots\]
Hom複体は
\[0\to \mr{Hom}(\cyc{2},\mb{Q}/\mb{Z}) \xrightarrow{2倍} \mr{Hom}(\mb{Z}/2\mb{Z},\mb{Q}/\mb{Z})\to 0 \to 0\to \dots\]
この2倍写像は$(a\mapsto f(a)) \mapsto (a \mapsto 2f(a))$である. 0,1次の部分のみでらうから, $n>1$なら$\mr{Ext}^{n}_{\mb{Z}}(\cyc{2},\cyc{2}) = 0$である. $\mr{Hom}(\cyc{2},\mb{Q}/\mb{Z})$は$1$の行き先で決定され, それは$0+\mb{Z}, \dfrac{1}{2}\mb{Z}$しかなく, $\mr{Hom}(\cyc{2},\mb{Q}/\mb{Z})\cong \cyc{2}$である(これは$\mr{Ext}^{0}_{\mb{Z}}$でもある). この同型において2倍写像は0写像になり, 像は0である. よって$\mr{Ext}^{1}_{\mb{Z}}(\cyc{2},\cyc{2}) = (\cyc{2})/0$.
\end{egans}




\subsection{演習問題}
\subsubsection{Level A}
\begin{prob}
アーベル群と$\mb{Z}$加群は同一視できることを示せ. すなわち, アーベル群から自然に$\mb{Z}$加群を与える方法が存在し, 逆に$\mb{Z}$加群からアーベル群を与える方法があり, それらが互いの逆対応であることを示せ.
\end{prob}

\begin{prob} $A$を\tf{局所環},
$M$を$A$加群, $x_1,\dots, x_n\in M$はとする. $x_i + \maf{m}M$が$k=A/\maf{m}$ベクトル空間$M/\maf{m}M$を$k$上生成するとき, $x_1,\dots, x_n$は$M$を$A$上生成することを示せ.
\end{prob}

\subsubsection{Level B}
\begin{prob} (\cite{liu}, Prop1.2.6)
$A$を環, $\mb{E}: 0\to M'\to M \to M'' \to 0$を$A$加群の完全列とする.
\ben
\item (\tf{復習})$M''$が射影的のとき, この完全列は分裂する. なぜか?
\item $M''$が\tf{平坦$A$加群であるとする.} このとき, 任意の$A$加群$N$について, $\mb{E}\tens_{A}N$はまた完全列であることを示したい. まず, $N$を自由加群の商$L/K$の形で表す. このとき, 次が可換である.
\[
\xymatrix{
 &M'\tens_{A}K\ar[r]\ar[d]& M\tens_{A}K\ar[r]\ar[d] &M''\tens_{A}K\ar[d]^{\alpha} \\
 0\ar[r]&M'\tens_{A}L\ar[r]\ar[d]^{\beta}&M\tens_{A}L\ar[r]\ar[d]& M''\tens_{A}L\\
  & M'\tens_{A}N \ar[r] & M\tens_{A}N
}
\]
これを用いて$M'\tens_{A}N \to M\tens_{A}N$の単射性を導け.
\een
\end{prob}

\begin{prob} (2020東工大午後2)
$K$は標数$p>0$の代数閉体とする. $r\geq 1$, $K$係数の$p^r$次多項式
\[f(X) = c_0X + c_1X^p + \dots + c_rX^{p^r} = \sum_{j=0}^{r} c_jX^{p^j}, c_0,c_r\neq 0\]
を固定する. $f$は写像$f:K\to K$; $x\mapsto f(x)$を定める. これを$k$回合成したものを$f^k$と表す. $A=\mb{F}_{p}[T]$の$K$への作用を, $a=\sum_{k=0}^{n} a_kT^k\in A$および$x\in K$に対し
\[ax = \sum_{k=0}^{n} a_kf^k(x)\]
により定める. これにより$K$は$A$加群となる. $a\in A$に対し$V_a = \set{x\in K\mid ax = 0}$とおく.
\ben
\item $a\neq 0$ならば$V_a$は有限生成$A$加群であることを示せ.
\item $a=T$のとき, $V_a$の$A$加群として構造を
\[A^{\oplus m}\oplus A/a_1A \oplus \dots \oplus A/a_lA\]
($m,l$は0以上の整数, $a_1,\dots, a_l\in A$)の形で求めよ.
\item $K$は$A$加群として有限生成ではないことを示せ.
\een
\end{prob}


\subsubsection{Level C}

\begin{prob}
$M$を\tf{局所環$A$上の有限生成平坦加群}とする. このとき$M$は\tf{自由加群}であることを示せ.
\end{prob}

\begin{prob}

\end{prob}




\subsubsection{Hint・Answer}
\begin{probans}
アーベル群$A$が与えられたとき, $a\in A$に対し$a+a+\dots + a$という$a$を$n$回足したものを$n\cdot a$と定めることで$\mb{Z}$加群の構造を入れることができる. 逆に, $\mb{Z}$加群は$\mb{Z}$の作用を忘れるとアーベル群である.
\end{probans}

\begin{probans}
$N=Ax_1+\dots + Ax_n$とする. $N\to M \to M/\maf{m}M$は全射なので$N+\maf{m}M = M$. ここから$M/N = 0$が, とある命題を用いて従う.
\end{probans}

\begin{probans}
平坦性より$\alpha$が単射である. $L\to N=L/K$が全射なので$\beta$も全射である. これだけを用いれば, diagram chasingで示せる.
\end{probans}

%===========================Level B==============================

\begin{probans}
(1): 「$K$が$A$加群になる」部分に$f(0)=0$が効いていることに注意. これにより, $b\in A$と$0\in K$に対し$b\cdot 0 = 0$が従う.
$x\in V_a$, $b\in A$なら, $K$が$A$加群なので$a(bx) = b(ax) = b\cdot 0 = 0$となる. 和で閉じていることは明らか. よって$V_a\subset K$は$A$加群となる. いま, $a_k\neq 0$である. $x\in V_a$であるための必要十分条件は$K$係数多項式$\sum_{k=0}^{n} a_kf^k(X)$の根になることであり, これは零多項式ではない. よってこれの根は有限個なので, 有限生成$A$加群である(その有限個の根で生成できる).\\
(2): $V_T$は$f(X)$の根の全体である. $b = \sum_{j=0}^{n} b_j T^j\in A$で$x\in V_T$を作用させると$bx = b_0x$であり,  特にイデアル$TA$の元に関する$V_T$への作用は0写像である. よって$V_T$は$A/TA \cong \mb{F}_{p}$ベクトル空間の構造と等価に与える. $V_T$は$p^r$個の元を持つので, $\mb{F}_{p}$上$r$個の元$x_1,\dots, x_r\in V_T$で生成される. つまり, $V_T$に自然に$\mb{F}_{p}$の作用を入れるとき$V_T = \mb{F}_p x_1 \oplus \dots \oplus \mb{F}_{p}x_r\cong \mb{F}_{p}^r$であり, この同型が$A$加群の同型$V_T \cong (A/TA)^{\oplus r}$を引き起こす. \\
(3): $K$が有限生成$A$加群であると仮定する. $A$はPIDだから単因子論により
\[K\cong A^{\oplus m} \oplus A/a_1A \oplus \dots \oplus A/a_l A\]
なる有限個の$a_i\in A$が存在する. $a_i$ ($i=1,\dots, l$)の次数の最大値を$N$とする. $A$には$N+1$次以上の既約多項式が存在するので, それをひとつ選び$F$とおく.

$F$の次数が$d$のとき, $A/FA \cong \mb{F}_{p^{d}}$である. (2)と同様に, $V_{F}$は$A/F A$加群の構造が定義され, $V_{F}$は$A/FA$のいくつかの直和に同型である. とくに, $K$のある元$x\neq 0$が$Fx = 0$を満たす. 単因子論による同型でこの$x$が同型で対応する元を$(a_1,a_2,\dots, a_m, r_1 + a_1A,r_2 + a_2A,\dots, r_l + a_lA)$と表すと, $Fa_i = 0\in A$, $Fr_j\in a_jA$が成り立つ.  $F$はノンゼロなので$a_i = 0$である. また, $b_j \in A$によって$Fr_j = a_j b_j$と表すとき, $F$の取り方から$F$と$a_j$は$A$で互いに素なので$b_j$は$F$で割り切れる. よって$r_j \in a_j A$であるから$x\neq 0$が対応する元は$0$になり, 同型であることに矛盾する. よって有限生成ではない. \qed
\end{probans}

\begin{probans}
$M/\maf{m}M$の$k=A/\maf{m}$上の基底を作る$x_1,\dots, x_n$を取れば, それが$M$を$A$上生成することは見た. それが一次独立であればよい. \par
必ずしも基底とは限らない一次独立系$\set{x_1+\maf{m},\dots, x_n+\maf{m}}\subset M/\maf{m}M$に対し, $\set{x_1,\dots,x_n}$が$A$上で一次独立であることを証明する. $\sum_{i} a_i x_i = 0$だとせよ. このとき$f:A^n \to A$を
\[f(b_1,b_2,\dots, b_n) = \sum a_i b_i\]
で定める. $0\to \ker{f}\hookrightarrow A^n \to \im{f} \to 0$に$\tens_{A}M$を施し$\ker{f}\tens_{A}M \to M^n \xrightarrow{f_M} M$という完全列を得る(ここで平坦を用いた). したがって$(x_1,\dots, x_n)$は$\ker{f} \tens_{A}M$の像に属す. 像と$\ker{f}\tens_{A}M$を同一視すれば, $(x_1,\dots, x_n)$は$\sum b_j \tens y_j$の形に表せる($b_j=(b_{j1},\dots, b_{jn})\in \ker{f}, y_j \in M$). $x_1,\dots, x_n$が$M/\maf{m}M$を生成するから,  すべての$b_{ij}$が$\maf{m}$に属すことはない. $b_{11}\in \maf{m}$としてよく, $A$は局所環なので$b_{11}\in A^{\times}$である. $\sum a_i x_i = 0$より$\sum_{i}a_ib_{i1} = 0$なので, $c_i = b_{i1}b_{11}^{-1}$として
\[a_1 + a_2c_2 + \dots + a_nc_n = 0\]
を得る. とくに$n=1$なら$a_1=0$なのでよい. $n\geq 2$なら, $\sum a_{i}x_i = 0$, $\sum_{i} a_ic_ix_1 = 0$ により
\[\sum_{i\geq 2} a_i(x_i - c_ix_1 ) = 0\]
である. $\set{x_i+\maf{m}M}$が$k$上一次独立ならば, 明らかに $\set{x_i - c_ix_1 + \maf{m}M}$もそうである. したがって帰納法を用いれば, 上から$a_i = 0$ ($i\geq 2$) が言え, $a_1=0$も従うので, $\set{x_1,\dots, x_n}$は一次独立になる. これにより, 基底$\set{x_1+\maf{m}M,\dots, x_n + \maf{m}M}$から引き戻して得た$M$の生成系$\set{x_1,\dots, x_n}$は基底でもあり, $A$は自由加群になる.
\lefba{ (\tf{振り返り}) $M/\maf{m}M$で自由なら$A$で自由であることを帰納法で示すのが難しい. $\sum a_ix_i = 0$として, 一次変換$f$を作って, 平坦性により$ker{f}\tens_{A}M\to M^n \to M$を得た. $(x_1,\dots, x_n)\in \ker{f}\tens_{A}M$だと思えるので, $\sum b_j\tens y_j$の形に書けば$b_j\in A^n$のある成分が$\maf{m}$に入らないので単元. それを$b_{11}$にした. $a_1 + c_2a_2 + \dots + c_n a_n = 0$という関係式を用いることで, $\sum a_ix_i$の$a_1x_1$を消すことができ, $n-1$個の生成系に対する関係式に帰着し, 帰納法が回る.

}

\end{probans}














%\section{($\spadesuit$) 次元}
%\subsection{整従属}
%\subsection{上昇定理・下降定理}
%\subsection{いろいろな例}
%\subsection{演習問題Part}


\newpage
\section{専門過去問特集}
\begin{prob} (H10数学II)
$\omega = \frac{-1+\sqrt{-3}}{2}$とする。
\ben
\item $(\mb{Z}[\omega])^{\times}$の位数を求めよ。
\item $(\mb{Z}[\omega]/(9))^{\times}$を巡回群の直和で表せ。
\een
\end{prob}

\begin{prob} (H14-II-1)

\end{prob}

\begin{prob} (H15-II-1)

\end{prob}

\begin{prob} (H15-II-2)

\end{prob}

\begin{prob} (H17-II-1)
$\mb{Z}[x_1,\cdots, x_n]$の極大イデアルは$n+1$個の元で生成されることを証明せよ。

\end{prob}

\begin{prob} (H18-II-2)

\end{prob}


